\documentclass[11pt]{amsart}
\usepackage{amsmath,amssymb,amsthm,mathtools,mathrsfs}
\usepackage[margin=1.08in]{geometry}
\usepackage{booktabs}
\usepackage{array}
\usepackage{longtable}
\usepackage{tabularx}
\usepackage[hypertexnames=false]{hyperref}
\hypersetup{hidelinks}
\usepackage[final,expansion=false]{microtype}
\setlength{\parskip}{4pt plus 1pt}
\emergencystretch=2.5em

\newtheorem{theorem}{Theorem}[section]
\newtheorem{proposition}[theorem]{Proposition}
\newtheorem{lemma}[theorem]{Lemma}
\newtheorem{corollary}[theorem]{Corollary}
\theoremstyle{definition}
\newtheorem{definition}[theorem]{Definition}
\theoremstyle{remark}
\newtheorem{remark}[theorem]{Remark}

\newcommand{\C}{\mathbb C}
\newcommand{\R}{\mathbb R}
\newcommand{\N}{\mathbb N}
\newcommand{\Z}{\mathbb Z}
\newcommand{\Q}{\mathbb Q}
\newcommand{\A}{\mathbb A}
\newcommand{\dd}{\,d}
\newcommand{\Tr}{\operatorname{Tr}}
\newcommand{\Rel}{\operatorname{Re}}
\newcommand{\Iml}{\operatorname{Im}}
\newcommand{\Res}{\operatorname*{Res}}
\newcommand{\supp}{\operatorname{supp}}
\newcommand{\vol}{\operatorname{vol}}
\newcommand{\rad}{\mathrm{rad}}
\newcommand{\DA}{\mathcal D_A}
\newcommand{\DAr}{\mathcal D_A^{\mathrm{rad}}}
\newcommand{\NA}{\mathcal N_A}
\newcommand{\CS}{C_S}
\newcommand{\cz}{\mathfrak c_0}
\newcommand{\BT}{B_T}
\newcolumntype{Y}{>{\raggedright\arraybackslash}X}

\begin{document}

\title[An explicit formula and a comparison theorem for a canonical trace]
{An Explicit Formula and a Comparison Theorem\\
for the Canonical Trace on a Correspondence Algebra\\
of the Idele Class Group}

\author{David Alejandro Trejo Pizzo}
\thanks{Independent Researcher. \texttt{dtrejopizzo@gmail.com}}

\date{\today}

\begin{abstract}
Let $C_S$ be the idele class group of $\Q$ attached to a finite set of places
$S$ containing the archimedean one, and let $T=\sum_{v\in S}W_v$ be the
canonical trace on the convolution algebra $C_c^\infty(C_S)$ constructed in
the companion papers.  We evaluate $T$ on radial test functions.  After
proving Weil's explicit formula in the normalization of $T$, by a contour
argument in which no contour crosses a zero, we show that for real
$F\in C_c^\infty(\R)$ the self-intersection of the radial class $\Gamma_F$ is
\[
  T_\infty(\Gamma_{F\star F^\dagger})
  =2\Phi_F(0)\Phi_F(1)-\sum_\rho\Phi_F(\rho)\Phi_F(1-\rho),
\]
where $\Phi_F$ is the Mellin transform of $F$ and $\rho$ runs over the zeros
of the completed zeta function.  The polar term $2\Phi_F(0)\Phi_F(1)$ is the
hyperbolic form of the Tate functionals, and it comes from the two simple
poles, of residue $1$, of the continued Euler-product symbol.  We also show
that the constant term of the trace is forced to vanish, that primitive
radial test functions suffice to detect a zero off the critical line, and
that the intersection form is nondegenerate on radial classes.
\end{abstract}

\maketitle

\tableofcontents

\section{Introduction}
\label{sec:intro}

Fix a finite set $S=\{\infty\}\cup S_f$ of places of $\Q$, with $S_f$ a finite
nonempty set of primes.  In \cite{P43,P44} we constructed, on the idele class
group
\[
  \CS\;\cong\;\R^\times\times\prod_{p\in S_f}\Z_p^\times ,
  \qquad |x|_S=|x_\infty|_\infty ,
\]
the commutative nuclear $\dagger$-algebra $\DA=C_c^\infty(\CS)$, with product
the convolution $\star$ for the multiplicative Haar measure $d^*x$ and
involution $F^\dagger(u)=\overline{F(u^{-1})}$, together with a trace
\[
  T(F)\;=\;\sum_{v\in S}W_v(F)
\]
whose local terms $W_p$, $W_\infty$ are the regularized Weil orbital
integrals of the local densities $\tau_v(x)=|x|_v^{1/2}/|1-x|_v$; the local
measures are normalized by $\vol_{d^*x}(p^k\Z_p^\times)=\log p$ and
$d^*x_\infty=dx/(2|x|)$.  The trace defines a Hermitian intersection form
$I_A(F,G)=T(F\star G^\dagger)$, and the two Tate functionals
$\lambda_\mp(F)=\int_{\CS}F(u)|u|_S^{\mp1/2}\,d^*u$ define a homomorphism
$M=(\lambda_-,\lambda_+)$ onto a hyperbolic plane $(\BT,\beta_T)$.  These
facts are recalled precisely in Section~\ref{sec:carrier}.

The subject of this paper is the value of $T$ on radial test functions, those
that depend only on $|w|_S$.  The answer is the explicit formula, and we prove
that formula here, in the normalization of $T$, rather than quote it.

Throughout, $\xi(s)=\tfrac12 s(s-1)\pi^{-s/2}\Gamma(s/2)\zeta(s)$ is the
completed zeta function, $\rho$ runs over the zeros of $\xi$ with
multiplicity, and for $F\in C_c^\infty(\R)$
\[
  \Phi_F(s)=\int_\R F(t)e^{(s-\frac12)t}\,dt ,\qquad
  \hat u(\tau)=\int_\R u(t)e^{-i\tau t}\,dt .
\]
For $h\in C_c^\infty(\R)$ the \emph{radial class} of $h$ is
$\Gamma_h(w)=h(-\log|w|_S)\in\DA$.

\medskip
\noindent\textbf{Theorem A} (Lemma~\ref{lem:radial-trace}).
\emph{For every $h\in C_c^\infty(\R)$,}
\[
  T(\Gamma_h)=\sum_{p\in S_f}\sum_{k\ge1}(\log p)\,p^{-k/2}
      \bigl[h(k\log p)+h(-k\log p)\bigr]
      +\frac1{2\pi}\int_\R \hat h(-\tau)\,\widehat W_\infty(0,\tau)\,d\tau ,
\]
\emph{the first sum being finite and the integral absolutely convergent.}

\medskip
\noindent\textbf{Theorem B} (Theorem~\ref{thm:EF}).
\emph{Let $F\in C_c^\infty(\R)$ be real with $\supp F\subseteq[-T,T]$ and put
$h=F\star F^\dagger$.  Then, the zero sum being taken in the symmetric order,}
\[
\begin{aligned}
  \sum_\rho\Phi_F(\rho)\Phi_F(1-\rho)
  ={}&2\Phi_F(0)\Phi_F(1)
   +\frac1{2\pi}\int_\R|\hat F(\tau)|^2
      \bigl[\Rel\psi(\tfrac14+\tfrac{i\tau}2)-\log\pi\bigr]d\tau\\
   &-\sum_{2\le n<e^{2T}}\frac{\Lambda(n)}{\sqrt n}
      \bigl[h(\log n)+h(-\log n)\bigr].
\end{aligned}
\]

\medskip
\noindent\textbf{Theorem C} (Theorem~\ref{thm:comparison}).
\emph{For real $F\in C_c^\infty(\R)$, with $T_\infty$ the trace on the
colimit over all $S$,}
\[
\begin{aligned}
  I_\infty(\Gamma_F,\Gamma_F)=T_\infty(\Gamma_{F\star F^\dagger})
  &=2\Phi_F(0)\Phi_F(1)-\sum_\rho\Phi_F(\rho)\Phi_F(1-\rho)\\
  &=\beta_T\bigl(M(\Gamma_F),M(\Gamma_F)\bigr)
     -\sum_\rho\Phi_F(\rho)\Phi_F(1-\rho).
\end{aligned}
\]

\medskip
\noindent\textbf{Theorem D} (Corollary~\ref{cor:c-unique}).
\emph{Among the traces $T^{(c)}=\sum_v W_v+c\,(\,\cdot\,)(1)$, $c\in\R$, the
identity of Theorem C holds only for $c=0$.}

\medskip
\noindent\textbf{Theorem E} (Theorem~\ref{thm:residues}).
\emph{Let $A(s)=\tfrac12\log\pi-\tfrac12\psi(s/2)-\zeta'/\zeta(s)$ be the
Euler-product symbol continued from $\Rel s>1$, and let
$A_S(s)=a(s)+\sum_{p\in S_f}f_p(s)$ be its truncation to $S$.  Then $A$ has
simple poles at $s=0$ and $s=1$, each of residue $1$, whereas
$\Res_{s=0}A_S=|S_f|+1$ and $\Res_{s=1}A_S=0$ for every finite $S$.
Evaluated against $g(s)=\Phi_F(s)\Phi_F(1-s)$, the two residues of $A$ give
exactly the polar term $2\Phi_F(0)\Phi_F(1)=\beta_T(M(\Gamma_F),M(\Gamma_F))$
of Theorem C.}

\medskip
\noindent\textbf{Theorem F} (Theorem~\ref{thm:sufficiency}).
\emph{If $\xi$ has a zero off the line $\Rel s=\tfrac12$, then there is a
nonzero real $F\in C_c^\infty(\R)$ with $\Phi_F(0)=\Phi_F(1)=0$ and}
\[
  \sum_\rho\Phi_F(\rho)\Phi_F(1-\rho)<0 ;
\]
\emph{$F$ may be taken even, and it may be taken odd.}

\medskip
\noindent\textbf{Theorem G} (Theorem~\ref{thm:radical}).
\emph{For every finite $S$ with $S_f\ne\emptyset$, the form $I_A$ is
nondegenerate on the radial classes.}

\medskip
Theorem B is Weil's explicit formula \cite{Weil1952} for a compactly
supported smooth test function; Barner \cite{Barner1981} gives it under
minimal regularity.  The identity is classical, and we prove it for two
reasons.  First, it is needed in exactly the normalization in which $T$ is
defined.  Second, the proof in Section~\ref{sec:EF} is arranged so that no
contour is shifted across a zero: the reflection $\xi'/\xi(1-s)=-\xi'/\xi(s)$
is applied before the logarithmic derivative is expanded, so that one works on
a single line $\Rel s=c>1$, and the shift to $\Rel s=\tfrac12$ is made
separately on each prime term, where the integrand $g(s)n^{-s}$ is entire.
Theorem C then reads the spectral side of the explicit formula as the trace
$T_\infty$ on radial classes, and Theorems D and E identify the constant and
the poles that make this identification exact.

The classical facts about the zeta function that we use are listed, with
sources, in Appendix~\ref{app:classical}.  Results of \cite{P43,P44} are
quoted as numbered propositions in Section~\ref{sec:carrier} and at a few
other points, each stated in full.

Section~\ref{sec:carrier} recalls the algebra, the trace, its symbol and the
Tate functionals.  Section~\ref{sec:radial} introduces the radial classes and
fixes conventions, and Section~\ref{sec:trace} proves Theorem A.
Sections~\ref{sec:normalization} and \ref{sec:arch} treat the normalization
of the local measures and the archimedean place, and
Section~\ref{sec:boundary} computes the Tate functionals on radial classes.
Section~\ref{sec:EF} proves Theorem B, Section~\ref{sec:comparison} Theorems C
and D, Section~\ref{sec:polar} Theorem E, Section~\ref{sec:sufficiency}
Theorem F and Section~\ref{sec:radicalsec} Theorem G.

\section{The trace and its symbol}
\label{sec:carrier}

Everything in this section is proved in \cite{P43}; we quote each statement
with its number there.

\subsection{The algebra}

\begin{proposition}[\cite{P43}, Lemma 2.2, Definition 2.4]
\label{prop:carrier}
Let $S=\{\infty\}\cup S_f$ with $S_f$ a finite nonempty set of primes, let
$\A_S^\times=\R^\times\times\prod_{p\in S_f}\Q_p^\times$ and let
$\Gamma_S\subset\A_S^\times$ be the group of $S$-units.  The valuation
homomorphism splits $\A_S^\times=\Gamma_S\times\ker\varphi$, so that the idele
class group
\[
  \CS:=\Gamma_S\backslash\A_S^\times\;\cong\;\ker\varphi\;\cong\;
  \R^\times\times\prod_{p\in S_f}\Z_p^\times
  \tag{2.1}
\]
is a locally compact abelian group, on which $|x|_S=|x_\infty|_\infty$.  Its
multiplicative Haar measure is normalized by
$\vol_{d^*x_p}(\Z_p^\times)=\log p$ at a finite place and by
$d^*x_\infty=dx/(2|x|)$ at the archimedean one.  The space
\[
  \DA:=C_c^\infty(\CS),\qquad
  (F\star G)(u)=\int_{\CS}F(v)G(uv^{-1})\,d^*v,\qquad
  F^\dagger(u)=\overline{F(u^{-1})},
  \tag{2.2}
\]
is a commutative nuclear LF $\dagger$-algebra.
\end{proposition}

For a prime $p\in S_f$ we write $j_p\colon\Q_p^\times\to\CS$ is the slice map that
places its argument at $p$ and $1$ elsewhere; the images
$j_p(p^k\Z_p^\times)$ are the \emph{valuation shells}.

\subsection{The trace}

\begin{proposition}[\cite{P43}, Definition 6.5 and Convention 6.10]
\label{prop:trace}
For $v\in S$ let $\tau_v(x)=|x|_v^{1/2}/|1-x|_v$ be the local Weil density.
There are continuous linear functionals $W_v$ on $\DA$, one for each $v\in S$,
defined by
\begin{align*}
  W_p(F)&=\sum_{k\ne0}(\log p)\,p^{-|k|/2}\!\!\int_{\Z_p^\times}\!\!
      (F\circ j_p)(p^ku)\,d^*u
     +(\log p)\!\int_{\Z_p^\times}\!\frac{(F\circ j_p)(u)-(F\circ j_p)(1)}{|1-u|_p}\,du,\\
  W_\infty(F)&=(\gamma+\log4\pi)\,b(0)
     +\int_0^\infty K(t)\bigl[b(t)+b(-t)-2b(0)e^{-t/2}\bigr]\,dt,
  \qquad K(t)=\frac{e^{-t/2}}{1-e^{-2t}},
\end{align*}
where $b$ denotes the profile of $F$ in the logarithmic coordinate
$t=\log|w|_S$.  The \emph{canonical trace} is
\[
  T(F):=\sum_{v\in S}W_v(F)+c\,F(1),\qquad c:=0 ,
  \tag{2.3}
\]
and more generally we write $T^{(c)}:=\sum_v W_v+c\,(\,\cdot\,)(1)$ for
$c\in\R$.  The \emph{intersection form} is
\[
  I_A(F,G):=T(F\star G^\dagger),
  \tag{2.4}
\]
Hermitian and jointly continuous on $\DA$, with trivial radical
$\mathcal N=0$ \emph{(\cite{P43}, Proposition 8.2 and Theorem 8.4)}.
\end{proposition}

\begin{proposition}[\cite{P43}, Theorem 6.6]
\label{prop:offidentity}
Let $\DA^{\mathrm{off}}\subset\DA$ be the subspace of test functions whose
support misses the identity class.  On $\DA^{\mathrm{off}}$ each $W_v$ is the
absolutely convergent integral $\int F(j_v(x))\tau_v(x)\,d^*x_v$; no
counterterm and no regularization occurs, and $W_v(F)=0$ for all but finitely
many $v$.
\end{proposition}

\subsection{The symbol}

Let $\widehat{\CS}$ be the dual group and $\cz$ its \emph{principal component},
the component of characters that are trivial on the compact part
$\prod_{p\in S_f}\Z_p^\times$ and on $\{\pm1\}$, parametrised by
$\tau\in\R$ through $w\mapsto|w|_S^{i\tau}$.

\begin{proposition}[\cite{P43}, Theorems 7.2 and 15.1(c)]
\label{prop:symbol}
There is a function $\widehat T_S$ on $\widehat{\CS}$, the \emph{symbol} of
$T$, with the following properties.
\begin{enumerate}
\item[(a)] On each component $\widehat T_S$ is real-analytic and real-valued,
and $\widehat T_S=\widehat W_\infty+\sum_{p\in S_f}\widehat W_p$ with
\[
  \widehat W_p(\cz,\tau)=2(\log p)\,\Rel\frac{\zeta_p}{1-\zeta_p},\qquad
  \zeta_p=p^{-1/2+i\tau},\qquad
  \widehat W_\infty(0,\tau)=\log\pi-\Rel\psi(\tfrac14+\tfrac{i\tau}2).
  \tag{2.5}
\]
\item[(b)] For every $c\in\R$ the function $\widehat T_S+c$ vanishes
identically on \emph{no} component of $\widehat{\CS}$, and its zero set on each
component is discrete.
\item[(c)] \emph{(Plancherel form.)}  For every $H\in\DA$,
\[
  T(H)=\int_{\widehat{\CS}}\widehat H(\chi)\bigl(\widehat T_S(\chi)+c\bigr)\,d\chi,
  \tag{2.6}
\]
the integral absolutely convergent, with $d\chi|_{\cz}=d\tau/2\pi$; and the
identity holds place by place, $W_v(H)=\int\widehat H\,\widehat W_v\,d\chi$
for every $v\in S$, each integral absolutely convergent.  No support hypothesis
is imposed on $H$.
\item[(d)] $\widehat T_S(\cz,\cdot)$ is even, and
$\widehat T_S(\cz,0)\ge\log\pi-\psi(\tfrac14)=\log\pi+\gamma+3\log2+\pi/2
\approx5.3722$.
\end{enumerate}
\end{proposition}

\subsection{The Tate functionals}

\begin{proposition}[\cite{P43}, Theorems 7.1 and 8.8]
\label{prop:tate}
The two \emph{Tate functionals}
\[
  \lambda_\mp(F):=\int_{\CS}F(u)|u|_S^{\mp1/2}\,d^*u
  \tag{2.7}
\]
are continuous on $\DA$, and $M:=(\lambda_-,\lambda_+)\colon\DA\to\BT:=\C^2$ is
a surjective continuous homomorphism of algebras satisfying
$M(F^\dagger)=M(F)^\sharp$, where $(z_-,z_+)^\sharp=(\bar z_+,\bar z_-)$.  The
target carries the hyperbolic Hermitian form
\[
  \beta_T(z,w)=z_-\overline{w_+}+z_+\overline{w_-},
  \tag{2.8}
\]
with matrix $\begin{psmallmatrix}0&1\\1&0\end{psmallmatrix}$ in the basis
$e_-,e_+$.
\end{proposition}

\begin{proposition}[\cite{P43}, Theorems 8.9(a), 8.11, 8.13--8.17]
\label{prop:rulings}
There is no $\DA$-module splitting of the Tate boundary inside $\DA$; the two
isotropic lines $\C e_\pm$ are realized only in the completion
$\mathcal N_A^{\mathrm{dist}}=\overline{J(\DA)}\subseteq\DA'\oplus\BT$ of
$\DA$ along $J(F)=(I_A(F,\cdot),M(F))$, as limits
$e_\pm^{\mathrm{dist}}=(0,e_\pm)$ of spread Tate characters.  The plane they span
is exactly the scalar radical of the completed form.
\end{proposition}

\subsection{Charts and the colimit}

\begin{proposition}[\cite{P43}, Theorem 11.5]
\label{prop:colimit}
The chart algebras $\DA(S)$ form a directed system under the maps
$J_{ST}\colon\DA(S)\to\DA(T)$ for $S\subseteq T$, with colimit
$\DA^\infty=\operatorname*{colim}_S\DA(S)$.  There is a chart-invariant trace
$T_\infty$ on $\DA^\infty$ with the following exactness property: if
$R(H)$ denotes the radius of the support of $H$ in the logarithmic coordinate,
then $W_\ell(J_{ST}H)=0$ for every prime $\ell\notin S_f$ with
$\log\ell>R(H)$.  Consequently $T_S(H)=T_\infty(H)$ whenever
$R(H)<\log\ell_{\mathrm{next}}(S)$, and outside that window the two differ by
the finite sum of the omitted local terms.  We write
$I_\infty(F,G):=T_\infty(F\star G^\dagger)$.
\end{proposition}

\begin{remark}
\label{rem:C1}
Every statement below in which the prime sum runs over all primes is a
statement about $T_\infty$ on $\DA^\infty$.  This does not apply to the
continued Euler-product symbol of Section~\ref{sec:polar}, which is an
analytic continuation and not a limit of chart symbols; see
Corollary~\ref{cor:noncommute}.
\end{remark}

Finally we recall that the form $I_A$ has no sign.

\begin{proposition}[\cite{P43}, Theorem 15.1]
\label{prop:indefinite}
With $c=0$, the form $I_A$ on $\DA$ is indefinite with infinite index on both
sides, and the primitive part $\{F: M(F)=0\}$ has trivial $I_A$-orthogonal
complement.
\end{proposition}

\section{Radial classes}
\label{sec:radial}

\begin{definition}[\cite{P44}, Definition 3.11]
\label{def:radial}
For $h\in C_c^\infty(\R)$ set
\[
  \Gamma_h(w):=h\bigl(-\log|w|_S\bigr),\qquad w\in\CS ,
\]
and
\[
  \DAr:=\{\Gamma_h:h\in C_c^\infty(\R)\}\subset\DA .
  \tag{3.1}
\]
Elements of $\DAr$ are called \emph{radial classes}.  We write
$R(\Gamma_h)=\sup\{|t|:t\in\supp h\}$.
\end{definition}

Each $\Gamma_h$ lies in $C_c^\infty(\CS)$: it is smooth because
$|\cdot|_S$ is, and compactly supported because the factor
$\prod_{p\in S_f}\Z_p^\times\times\{\pm1\}$ is compact and $h$ has compact
support.

\begin{proposition}[\cite{P44}, Lemma 3.12]
\label{prop:radalg}
$\DAr$ is a $\dagger$-stable subalgebra of $\DA$ and $h\mapsto\Gamma_h$ is an
isomorphism of $\dagger$-algebras from $(C_c^\infty(\R),*)$ onto it:
\[
  \Gamma_h\star\Gamma_g=\Gamma_{h*g},\qquad
  \Gamma_h^\dagger=\Gamma_{h^\vee},\quad h^\vee(t)=\overline{h(-t)},\qquad
  \lambda_\mp(\Gamma_h)=\int_\R h(t)e^{\pm t/2}\,dt ,
  \tag{3.2}
\]
and $0\notin\supp h$ implies $\Gamma_h\in\DA^{\mathrm{off}}$.  Injectivity is
immediate from \emph{(2.1)}, since $|\cdot|_S$ is surjective onto $\R_{>0}$.
\end{proposition}

\begin{remark}
\label{rem:typing}
The Haar measures of the individual shells $j_p(p^k\Z_p^\times)$, and the
Frobenius measures of \cite{P44}, are not elements of $\DA$; by \cite{P44}, Theorems 4.2, 4.7 and 6.11, they are realized only in the completion
$\mathcal N_A^{\mathrm{dist}}$, as limits in the strong dual, and the numbers
$\Lambda(n)n^{-1/2}$ are the atom masses of the unique representing measure of
the finite part of $T$.  All test functions in this paper are elements of
$\DAr$, except in Proposition~\ref{prop:live}.
\end{remark}

\subsection{Coordinates}

Three coordinates occur.  On $\R_{>0}^\times$ one may use the multiplicative
variable $x$; on $\R$ the additive variable $t=\log x$; on $\CS$ the modulus
coordinate $t=\log|w|_S$.  Under the disintegration $d^*w=dt\otimes\mu$, with
$\mu$ the normalized (probability) Haar measure of the compact part
\cite[(10)]{P43}, the modulus coordinate identifies the
radial classes with $C_c^\infty(\R)$.  The unnormalized multiplicative
involution $f^\dagger(x)=x^{-1}\overline{f(x^{-1})}$ is carried to
$F^\dagger(u)=\overline{F(u^{-1})}$ by $F\mapsto|\cdot|_S^{1/2}F$, which in the
additive coordinate is $f\mapsto e^{t/2}f(e^t)$; so the half-density is the
change of variables between the two conventions.

\begin{lemma}
\label{lem:profile}
For $h\in C_c^\infty(\R)$ the profile of $\Gamma_h$ in the coordinate
$t=\log|w|_S$ is $t\mapsto h(-t)$.  Consequently
\[
  \widehat{\Gamma_h}(\cz,\tau)=\hat h(-\tau),
  \tag{3.3}
\]
and $\widehat{\Gamma_h}$ vanishes on every non-principal component of
$\widehat{\CS}$.
\end{lemma}

\begin{proof}
The profile is Definition~\ref{def:radial}.  The function $\Gamma_h$ is
constant on the compact coordinates and even in the sign of $w_\infty$, so by
orthogonality of characters its Fourier transform is carried by the component
$\cz$ on which all the compact characters and the sign character are trivial.
On that component the character is $w\mapsto|w|_S^{i\tau}$, so
\[
  \widehat{\Gamma_h}(\cz,\tau)=\int_{\CS}\Gamma_h(w)|w|_S^{-i\tau}\,d^*w
  =\int_\R h(-t)e^{-i\tau t}\,dt=\hat h(-\tau).\qedhere
\]
\end{proof}

Since every symbol occurring below is even in $\tau$ on $\cz$
(Proposition~\ref{prop:symbol}(d)), the sign in (3.3) never changes a value of
$I_A$; see Remark~\ref{rem:swap}.

\subsection{Conventions}

We fix
\begin{equation}
\begin{aligned}
  \hat u(\tau)&=\int_\R u(t)e^{-i\tau t}\,dt,
  &u(t)&=\frac1{2\pi}\int_\R\hat u(\tau)e^{i\tau t}\,d\tau,\\
  \Phi_F(s)&=\int_\R F(t)e^{(s-\frac12)t}\,dt,
  &\Phi_F(\tfrac12+i\tau)&=\hat F(-\tau),
\end{aligned}
\tag{3.4}
\end{equation}
so that, by (3.2),
\begin{equation}
  \lambda_-(\Gamma_F)=\Phi_F(1),\qquad \lambda_+(\Gamma_F)=\Phi_F(0),
  \qquad d\chi|_{\cz}=\frac{d\tau}{2\pi}.
\tag{3.5}
\end{equation}

For later reference we record the archimedean constants, all of them
elementary consequences of Proposition~\ref{prop:symbol} and of the digamma facts (E4) of
Appendix~\ref{app:classical}:
\begin{equation}
\begin{aligned}
  \widehat W_\infty(0,\tau)&=\log\pi-\Rel\psi(\tfrac14+\tfrac{i\tau}2),\\
  \widehat W_\infty(0,0)&=\log\pi-\psi(\tfrac14)
    =\log\pi+\gamma+3\log2+\tfrac\pi2=5.372184\ldots,\\
  \gamma+\log4\pi+\psi(\tfrac12)&=\log\pi
   \qquad\text{(since }\psi(\tfrac12)=-\gamma-2\log2\text{)},\\
  K(t)&=\frac{e^{-t/2}}{1-e^{-2t}}
    =\frac14\Bigl[\frac1{\sinh(t/2)}+\frac1{\cosh(t/2)}\Bigr],
  \qquad K(\log n)=\frac{n^{3/2}}{n^2-1}.
\end{aligned}
\tag{3.6}
\end{equation}
The third line is the reconciliation of the counterterm $\gamma+\log4\pi$ of
$W_\infty$ with the constant $\log\pi$ of its symbol; it is used in
Theorem~\ref{thm:archsymbol}.

\section{The radial trace}
\label{sec:trace}

\subsection{Test functions vanishing near the identity}

\begin{proposition}[\cite{P44}, Theorem 4.1]
\label{prop:primepacket}
Let $h\in C_c^\infty(\R)$ with $0\notin\supp h$, and let $c=0$.  Then
$T(\Gamma_h)$ is an absolutely convergent finite sum plus a convergent
integral, with no regularization and no counterterm, and
\begin{align}
  T(\Gamma_h)
  &=\sum_{p\in S_f}\sum_{k\ge1}(\log p)\,p^{-k/2}
     \bigl[h(k\log p)+h(-k\log p)\bigr]+\int_\R K(|t|)h(t)\,dt
  \tag{4.1}\\
  &=\sum_{n=p^k,\;p\in S_f}\Lambda(n)n^{-1/2}
     \bigl[h(\log n)+h(-\log n)\bigr]+\int_\R K(|t|)h(t)\,dt .
  \tag{4.2}
\end{align}
The proof has two steps: \emph{(i)} on a radial class the regularized unit
term of $W_p$ vanishes identically and pointwise, because
$|j_p(u)|_S=1$ for $u\in\Z_p^\times$, so that $\Gamma_h\circ j_p$ is constant
on $\Z_p^\times$; \emph{(ii)} what remains is the shell part, computed from
$\tau_p\equiv p^{-|k|/2}$ on $p^k\Z_p^\times$ \emph{(\cite{P43}, Lemma 6.2(i))}
and $\vol_{d^*x}(p^k\Z_p^\times)=\log p$.
\end{proposition}

Two remarks on (4.2).  First, only prime powers occur, and this is a fact
about $\CS$ rather than an input from number theory: the finite part of $T(\Gamma_h)$ is a sum over the valuation shells
$p^k\Z_p^\times$, whose moduli are exactly the prime powers, and no single
place has a shell of modulus $1/n$ for composite $n$.  (At the level of
objects the corresponding statement is \cite{P44}, Theorem 6.8: for
composite $n$ the support of the Frobenius class meets no shell at all, whereas
for $n=p^k$ it is the whole shell $j_p(p^k\Z_p^\times)$.  That is a statement
about supports, prior to any trace.)
Second, in the atom mass $\Lambda(p^k)p^{-k/2}$ \emph{both} factors are local
data of the single finite place $p$: $p^{-k/2}$ is the constant value of the
local density $\tau_p$ on the shell, and $\log p$ is the shell's
$d^*x$-volume.  The archimedean place contributes no atom at all; its
contribution is the absolutely continuous density $K(|t|)\,dt$.

\subsection{The general case}

Proposition~\ref{prop:primepacket} is stated for $0\notin\supp h$, and the
excluded case is the one needed below: every application in this
paper has $h=F\star F^\dagger$, so that $h(0)=\|F\|_{L^2}^2>0$ and
$0\in\supp h$ always.  Nor can the identity be obtained from (4.1) by
continuity, because the naive archimedean integral $\int K(|t|)h(t)\,dt$
diverges at $t=0$, since $K(t)\sim1/(2t)$.  The correct statement replaces the
naive integral by the Plancherel form of the archimedean term.

\begin{lemma}[Theorem A]
\label{lem:radial-trace}
Let $h\in C_c^\infty(\R)$ be arbitrary and let $c=0$.  Then $\Gamma_h\in\DA$,
$T(\Gamma_h)$ is defined, and
\begin{equation}
  T(\Gamma_h)=\sum_{p\in S_f}\sum_{k\ge1}(\log p)\,p^{-k/2}
     \bigl[h(k\log p)+h(-k\log p)\bigr]
   +\frac1{2\pi}\int_\R\hat h(-\tau)\,\widehat W_\infty(0,\tau)\,d\tau ,
\tag{4.3}
\end{equation}
the finite part a finite sum and the archimedean part an absolutely convergent
integral.  When $0\notin\supp h$, the second term equals $\int_\R
K(|t|)h(t)\,dt$ and \emph{(4.3)} reduces to \emph{(4.1)}.  On the colimit, with
$T_\infty$ in place of $T$, the restriction $p\in S_f$ disappears and the first
sum runs over all primes (still a finite sum, by compactness of $\supp h$).
\end{lemma}

\begin{proof}
We treat the two place types separately; neither argument uses a support
hypothesis.

\emph{The finite places.}  By Proposition~\ref{prop:trace},
$T(\Gamma_h)=\sum_{v\in S}W_v(\Gamma_h)$ with $c=0$, valid on all of $\DA$.  At
a finite place $W_p$ is the sum of a shell integral and the regularized unit
term
\[
  (\log p)\int_{\Z_p^\times}\frac{a(u)-a(1)}{|1-u|_p}\,du ,
  \qquad a:=\Gamma_h\circ j_p .
\]
For a radial class $a$ is \emph{constant} on $\Z_p^\times$: for a unit $u$ the
idele $j_p(u)$ has $|j_p(u)|_S=1$, hence $a(u)=h(0)=a(1)$ for every
$u\in\Z_p^\times$.  The unit term therefore vanishes identically.  What
remains is the shell part.  By \cite{P43}, Lemma 6.2(i), the local
density $\tau_p$ is
constant on the shell $p^k\Z_p^\times$, with value $p^{-|k|/2}$; and
$\Gamma_h\circ j_p$ is constant on that shell, with value $h(k\log p)$, because
$|j_p(p^ku)|_S=p^{-k}$ for $u\in\Z_p^\times$.  Since
$\vol_{d^*x}(p^k\Z_p^\times)=\log p$, the shell part of $W_p(\Gamma_h)$ is
\[
  \sum_{k\ne0}(\log p)\,p^{-|k|/2}\,h(k\log p)
  =\sum_{k\ge1}(\log p)\,p^{-k/2}\bigl[h(k\log p)+h(-k\log p)\bigr],
\]
the finite sum displayed in (4.3); it is finite because $\supp h$ is compact,
so that only the finitely many $k$ with $|k|\log p\le R(\Gamma_h)$ contribute.

\emph{The archimedean place.}  Here we use the Plancherel form of $W_\infty$
instead of its definition.  By Proposition~\ref{prop:symbol}(c) the
identity $W_v(H)=\int_{\widehat{\CS}}\widehat H\,\widehat W_v\,d\chi$ holds for
every $v\in S$ and every $H\in\DA$, each integral absolutely convergent and no
support hypothesis imposed.  Apply this with $H=\Gamma_h$.  By
Lemma~\ref{lem:profile} the transform $\widehat{\Gamma_h}$ is carried by $\cz$
and equals $\hat h(-\tau)$ there, and $d\chi|_{\cz}=d\tau/2\pi$; hence
\[
  W_\infty(\Gamma_h)=\frac1{2\pi}\int_\R\hat h(-\tau)\,\widehat
  W_\infty(0,\tau)\,d\tau ,
\]
absolutely convergent.  Adding the finite-place terms computed above gives
(4.3).  (Equivalently: the counterterm $(\gamma+\log4\pi)b(0)$ and the
subtraction $-2b(0)e^{-t/2}$ in the definition of $W_\infty$ are exactly what
makes $W_\infty$ a continuous functional on all of $\DA$ with symbol
$\widehat W_\infty$; the divergence of the naive $\int K(|t|)h(t)\,dt$ at $t=0$
is what they remove, and (4.3) is the regularized value.)

\emph{Consistency.}  If $0\notin\supp h$ then $b(0)=0$, both the counterterm
and the subtraction vanish (Proposition~\ref{prop:offidentity}), and the
archimedean term is
$\int_0^\infty K(t)[h(-t)+h(t)]\,dt=\int_\R K(|t|)h(t)\,dt$, so (4.3) reduces
to (4.1).

\emph{The colimit statement} is Proposition~\ref{prop:colimit}: the local terms
$W_\ell$ for $\ell\notin S_f$ with $\log\ell>R(\Gamma_h)$ vanish identically,
so the sum over all primes is finite and equals the chart sum for every chart
containing the finitely many primes $\ell\le e^{R(\Gamma_h)}$.
\end{proof}

\section{The normalization of the local measures}
\label{sec:normalization}

The normalization $\vol_{d^*x}(p^k\Z_p^\times)=\log p$ used above is forced, up to one global scale, by the
requirement that the local Tate residue be the same at every place.

\begin{proposition}[\cite{P44}, Theorem 11.1 and Corollary 11.2]
\label{prop:normalization}
Normalize the local multiplicative measures by
$\vol_{d^*x_p}(\Z_p^\times)=c_p>0$ and $d^*x_\infty=dx/(2|x|)$, all other
conventions unchanged, and define the \emph{local Tate residue} of a place $v$
as $\Res_{s=0}\zeta_v(s)$ with
$\zeta_v(s)=\int_{|x|_v\le1}|x|_v^s\,d^*x_v$.  Then
\[
  \zeta_p(s)=\frac{c_p}{1-p^{-s}},\qquad
  \Res_{s=0}\zeta_p=\frac{c_p}{\log p},\qquad
  \zeta_\infty(s)=\frac1s,\qquad \Res_{s=0}\zeta_\infty=1 ,
\]
so the local Tate residues agree at all places if and only if $c_p=\log p$ for
every finite $p$.  The only remaining freedom is a single global positive
scale, which multiplies $I_A$ by a positive constant and changes no sign, no
vanishing and no ratio.
\end{proposition}

\begin{corollary}
\label{cor:threenumbers}
Under $c_p=\log p$ the following three numbers are equal to $\log p$:
\emph{(i)} $\vol_{d^*x}(p^k\Z_p^\times)$, by
Proposition~\ref{prop:normalization}; \emph{(ii)} $\Lambda(p^k)$, the atom mass
of the finite part of $T$, by Proposition~\ref{prop:primepacket}; and
\emph{(iii)} $\log\#(\Z/p)$.  The equality of \emph{(i)} and \emph{(ii)} is a
theorem; the equality with \emph{(iii)} is a coincidence of values, and no map
relating \emph{(iii)} to \emph{(i)} is asserted.
\end{corollary}

\begin{proof}
(i) and (ii) are the two cited statements; (iii) is clear.
\end{proof}

\section{The archimedean place}
\label{sec:arch}

The archimedean term of $T$ is the regularized integral of the local density
$\tau_\infty$ against the profile.  This section identifies its symbol, records
the spectral reading of its kernel, and computes the archimedean Tate residue,
which is the archimedean half of the residue comparison of
Section~\ref{sec:polar}.

\subsection{The kernel as a heat trace}

\begin{definition}
\label{def:NGamma}
Let $N_\Gamma$ denote the self-adjoint operator with simple spectrum
$\{2m+\tfrac12: m\ge0\}$.  It is \emph{defined} by its spectrum; no independent
construction of it is used or claimed.
\end{definition}

\begin{proposition}
\label{prop:heattrace}
For $t>0$,
\[
  K(t)=\frac{e^{-t/2}}{1-e^{-2t}}=\sum_{m\ge0}e^{-(2m+\frac12)t}
  =\Tr e^{-tN_\Gamma}.
\]
\end{proposition}

\begin{proof}
Geometric series, convergent for $t>0$.
\end{proof}

\begin{proposition}[the zeta determinant]
\label{prop:zetadet}
For $\Rel s>0$,
\begin{equation}
  {\det}_\zeta\bigl(N_\Gamma+s-\tfrac12\bigr)
  =\frac{\sqrt{2\pi}\,2^{\frac12-\frac s2}}{\Gamma(s/2)} .
\tag{6.1}
\end{equation}
\end{proposition}

\begin{proof}
The operator $N_\Gamma+s-\tfrac12$ has simple spectrum $\{2m+s: m\ge0\}$, so
its spectral zeta function is
\[
  \zeta_{N_\Gamma,s}(w)=\sum_{m\ge0}(2m+s)^{-w}=2^{-w}\zeta_H(w,s/2),
\]
with $\zeta_H$ the Hurwitz zeta function, regular at $w=0$.  Differentiating,
\[
  \zeta_{N_\Gamma,s}'(0)=-(\log2)\zeta_H(0,a)+\zeta_H'(0,a),\qquad a=s/2 .
\]
Insert $\zeta_H(0,a)=\tfrac12-a$ and Lerch's formula
$\zeta_H'(0,a)=\log\Gamma(a)-\tfrac12\log2\pi$ (fact (E7) of
Appendix~\ref{app:classical}), and exponentiate
${\det}_\zeta=\exp(-\zeta_{N_\Gamma,s}'(0))$:
\[
  {\det}_\zeta(N_\Gamma+s-\tfrac12)
  =\exp\Bigl[(\log2)(\tfrac12-\tfrac s2)-\log\Gamma(\tfrac s2)+\tfrac12\log2\pi\Bigr]
  =\frac{\sqrt{2\pi}\,2^{\frac12-\frac s2}}{\Gamma(s/2)} .\qedhere
\]
\end{proof}

\begin{remark}
\label{rem:scope}
Formula (6.1) gives $\Gamma(s/2)$ only up to the scalar
$\sqrt{2\pi}\,2^{\frac12-\frac s2}$, and the factor $\pi^{-s/2}$ of the
completed zeta function is \emph{not} produced by $N_\Gamma$: reaching
$\pi^{-s/2}\Gamma(s/2)$ requires the extra scalar $\sqrt\pi\,(2\pi)^{-s/2}$.
The $\log\pi$ occurring in (6.2) below comes from the independently computed
finite part of $W_\infty$, not from $N_\Gamma$.  None of Sections~\ref{sec:EF}--\ref{sec:radicalsec} uses
Proposition~\ref{prop:heattrace} or Proposition~\ref{prop:zetadet}.
\end{remark}

\subsection{The archimedean symbol}

\begin{theorem}
\label{thm:archsymbol}
With $W_\infty$ as in Proposition~\emph{\ref{prop:trace}},
\begin{equation}
  \widehat W_\infty(0,\tau)=\log\pi-\Rel\psi(\tfrac14+\tfrac{i\tau}2),
\tag{6.2}
\end{equation}
and the counterterm constant reconciles as
$\gamma+\log4\pi+\psi(\tfrac12)=\log\pi$.
\end{theorem}

\begin{proof}
Identity (6.2) is Proposition~\ref{prop:symbol}(a).  For the constant, insert
$\psi(\tfrac12)=-\gamma-2\log2$ (fact (E4)):
$\gamma+\log4\pi-\gamma-2\log2=\log4\pi-\log4=\log\pi$.
\end{proof}

\subsection{The archimedean Tate residue}

\begin{theorem}
\label{thm:archresidue}
Put $s=\tfrac12+i\tau$, so that $\tau=i/2$ corresponds to $s=0$ and
$\tau=-i/2$ to $s=1$, and $ds=i\,d\tau$.  Set
\[
  a(s):=\tfrac12\log\pi-\tfrac12\psi(s/2).
\]
Then:
\begin{enumerate}
\item[(i)] $\widehat W_\infty(0,\tau)=a(s)+a(1-s)$ as meromorphic functions;
\item[(ii)] $a$ has a simple pole at $s=0$ with residue exactly $1$ and no
other pole in $\Rel s>-2$; hence $\widehat W_\infty(0,\cdot)$ has a simple pole
at $\tau=i/2$ with residue $-i$ and at $\tau=-i/2$ with residue $+i$;
\item[(iii)] $\Res_{s=0}a(s)=1=\Res_{s=0}\zeta_\infty(s)$, where
$\zeta_\infty(s)=\int_{|x|\le1}|x|^s\,dx/(2|x|)=1/s$.  The same holds at every
finite place: with $f_p(s):=(\log p)p^{-s}/(1-p^{-s})$ one has
$\Res_{s=0}f_p=1=\Res_{s=0}\zeta_p(s)$, the latter precisely when
$c_p=\log p$.
\end{enumerate}
\end{theorem}

\begin{proof}
(i)  For real $\tau$ and $s=\tfrac12+i\tau$,
\[
  a(s)+a(1-s)=\log\pi-\tfrac12\bigl[\psi(\tfrac14+\tfrac{i\tau}2)
  +\psi(\tfrac14-\tfrac{i\tau}2)\bigr]
  =\log\pi-\Rel\psi(\tfrac14+\tfrac{i\tau}2),
\]
because $\psi(\bar z)=\overline{\psi(z)}$.  This is (6.2).  Both sides are
meromorphic in $s$ and agree on a line, hence everywhere.

(ii)  By (E4), $\psi(z)=-1/z+O(1)$ as $z\to0$, so
$a(s)=-\tfrac12\cdot(-2/s)+O(1)=1/s+O(1)$; the next pole of $\psi(s/2)$ is at
$s=-2$.  For the transport to the $\tau$-variable: if $G$ has residue $r$ at
$s_0$, then as a function of $\tau$ with $s=\tfrac12+i\tau$ its residue at
$\tau_0=(s_0-\tfrac12)/i$ is $r\,(ds/d\tau)^{-1}=r/i=-ir$.  With $r=1$ at
$s=0$ this is $-i$ at $\tau=i/2$; the pole at $s=1$ comes from $a(1-s)$, whose
residue there is $-1$ in $s$, hence $+i$ in $\tau$ at $\tau=-i/2$.

(iii)  $\int_{|x|\le1}|x|^s\,dx/(2|x|)=2\cdot\tfrac12\int_0^1x^{s-1}dx=1/s$,
of residue $1$.  At a finite place,
$\zeta_p(s)=c_p/(1-p^{-s})$ has residue $c_p/\log p$ at $s=0$ by
Proposition~\ref{prop:normalization}, while
$1-p^{-s}=s\log p+O(s^2)$ gives $f_p(s)=1/s+O(1)$, i.e.\ $\Res_{s=0}f_p=1$.
\end{proof}

\begin{remark}
\label{rem:onepole}
Theorem~\ref{thm:archresidue} says that the pole of the \emph{symbol} of $W_v$
at $\tau=\mp i/2$ is, place by place and under the dictionary
$s=\tfrac12+i\tau$, the pole at $s=0,1$ of the \emph{local Tate zeta integral},
with the same residue at every place --- and that the requirement ``the same at
every place'' is exactly what forces $\vol(\Z_p^\times)=\log p$
(Proposition~\ref{prop:normalization}).  What happens to these residues when
one passes from a chart to the full Euler product is the subject of
Section~\ref{sec:polar}.
\end{remark}

\section{The Tate functionals on radial classes}
\label{sec:boundary}

\begin{theorem}
\label{thm:polarpacket}
Let $(\BT,\beta_T)$ and $M=(\lambda_-,\lambda_+)$ be as in
Proposition~\emph{\ref{prop:tate}}, and let $F\in C_c^\infty(\R)$ be real.
Then:
\begin{enumerate}
\item[(i)] $M(\Gamma_F)=\bigl(\Phi_F(1),\Phi_F(0)\bigr)$ and
\[
  \beta_T\bigl(M(\Gamma_F),M(\Gamma_F)\bigr)=2\Phi_F(0)\Phi_F(1)
  =\Phi_h(0)+\Phi_h(1),\qquad h=F\star F^\dagger ;
\]
\item[(ii)] $M(\Gamma_F)=0$ if and only if $\Phi_F(0)=\Phi_F(1)=0$; we call
such an $F$, and the class $\Gamma_F$, \emph{primitive};
\item[(iii)] the involution acts on $\BT$ by $\sharp$, exchanging the two
isotropic lines.
\end{enumerate}
\end{theorem}

\begin{proof}
(i)  By (3.2), $\lambda_\mp(\Gamma_F)=\int F(t)e^{\pm t/2}dt$, so
$\lambda_-(\Gamma_F)=\Phi_F(1)$ and $\lambda_+(\Gamma_F)=\Phi_F(0)$, which is
(3.5).  For real $F$ both are real, so
\[
  \beta_T(M,M)=\lambda_-\overline{\lambda_+}+\lambda_+\overline{\lambda_-}
  =2\Phi_F(0)\Phi_F(1).
\]
For the last equality, $\Phi_{F^\dagger}(s)=\Phi_F(1-s)$ for real $F$
(Lemma~\ref{lem:testfn} below), so
$\Phi_h(s)=\Phi_F(s)\Phi_F(1-s)$ and $\Phi_h(0)=\Phi_F(0)\Phi_F(1)=\Phi_h(1)$.

(ii) is (3.5).  (iii) is $\lambda_\mp(F^\dagger)=\overline{\lambda_\pm(F)}$,
which is Proposition~\ref{prop:tate}.
\end{proof}

\begin{remark}[the swap]
\label{rem:swap}
Writing $\sigma$ for the transposition of the two coordinates, part (i) reads
$M(\Gamma_F)=\sigma\bigl(\Phi_F(0),\Phi_F(1)\bigr)$: the boundary value is the
\emph{swap} of the pair $(\Phi_F(0),\Phi_F(1))$, not that pair itself.  The
distinction decides which isotropic direction a given family of test classes
approximates, and it does not affect any quadratic statement, because $\beta_T$ is $\sigma$-invariant,
\[
  \beta_T(\sigma z,\sigma w)=(\sigma z)_-\overline{(\sigma w)_+}
  +(\sigma z)_+\overline{(\sigma w)_-}
  =z_+\overline{w_-}+z_-\overline{w_+}=\beta_T(z,w).
\]
With the normalization of Proposition~\ref{prop:tate}, the spread families of
\cite{P43}, Theorem 8.13, have $\lambda_+\to1$ and $\lambda_-\to0$, so they
approximate $e_+$; and since $\lambda_+(\Gamma_F)=\Phi_F(0)$, the direction
$e_+$ pairs with the value of $\Phi_F$ at $s=0$, that is, with the pole of
$\xi$'s trivial factor at $s=0$.
\end{remark}

\begin{proposition}[\cite{P43}, Theorems 8.9(a), 8.13, 8.15, 8.17 and
Definition 8.16]
\label{prop:live}
The isotropic directions $e_\pm$ are not classes of $\DA$: by
Proposition~\emph{\ref{prop:rulings}} there is no $\DA$-module splitting of the
Tate boundary inside $\DA$, so the boundary component of a test class must be
taken through $M$.  They exist in the completion
$\mathcal N_A^{\mathrm{dist}}$, as limits of spread Tate characters, and the
plane they span is exactly the scalar radical there.  In particular they are
not radial classes: they are limits of the spread families of
\cite{P43}, Definition 8.16, whose $L^1$-mass escapes to $|w|_S\to\infty$.
\end{proposition}

\begin{proof}
The three assertions are, in order, \cite{P43}, Theorem 8.9(a),
Theorems 8.13 and 8.15, and Definition 8.16 together with
Theorem 8.17.
\end{proof}

\section{The explicit formula}
\label{sec:EF}

This section proves Theorem B, using only the facts (E1)--(E5) of
Appendix~\ref{app:classical}.

\subsection{Test functions}

\begin{lemma}
\label{lem:testfn}
Let $F\in C_c^\infty(\R)$ be real with $\supp F\subseteq[-T,T]$, and let
$\Phi_F(s)=\int F(t)e^{(s-\frac12)t}dt$.  Then $\Phi_F$ is entire,
$\Phi_F(\bar s)=\overline{\Phi_F(s)}$, and for every integer $N\ge0$, with
$C_N:=\|F^{(N)}\|_{L^1}$,
\begin{equation}
  |\Phi_F(\sigma+i\tau)|\le C_N\,e^{T|\sigma-\frac12|}\,|s-\tfrac12|^{-N}
  \qquad(s=\sigma+i\tau\ne\tfrac12).
\tag{8.1}
\end{equation}
Put $g(s):=\Phi_F(s)\Phi_F(1-s)$.  Then $g$ is entire, $g(1-s)=g(s)$,
$g(\bar s)=\overline{g(s)}$, $g(\tfrac12+i\tau)=|\hat F(\tau)|^2\ge0$, and for
every $N$
\begin{equation}
  |g(\sigma+i\tau)|\le C_N^2\,e^{2T|\sigma-\frac12|}\,|s-\tfrac12|^{-2N}.
\tag{8.2}
\end{equation}
Moreover, with $h:=F\star F^\dagger$, the function $h$ is real and even,
supported in $[-2T,2T]$, satisfies $\hat h=|\hat F|^2$ and
$h(a)=\langle S_aF,F\rangle$ with $(S_aF)(t)=F(t+a)$, and $\Phi_h=g$.
\end{lemma}

\begin{proof}
Entirety is differentiation under the integral sign over a compact set, and
$\Phi_F(\bar s)=\overline{\Phi_F(s)}$ holds because $F$ is real.

For (8.1), integrate by parts $N$ times; all boundary terms vanish by compact
support, and
\[
  \Phi_F(s)=\Bigl(\tfrac12-s\Bigr)^{-N}\int_\R F^{(N)}(t)e^{(s-\frac12)t}\,dt
  \quad\text{up to a sign,}
\]
whence $|\Phi_F(s)|\le|s-\tfrac12|^{-N}\|F^{(N)}\|_{L^1}e^{T|\sigma-\frac12|}$,
because $|e^{(s-\frac12)t}|=e^{(\sigma-\frac12)t}\le e^{T|\sigma-\frac12|}$ on
$\supp F$.

$g$ is entire as a product of entire functions; $g(1-s)=g(s)$ is the
definition; $g(\bar s)=\overline{g(s)}$ follows from
$\Phi_F(\bar s)=\overline{\Phi_F(s)}$.  On the line $s=\tfrac12+i\tau$ one has
$\Phi_F(\tfrac12+i\tau)=\int F(t)e^{i\tau t}dt=\hat F(-\tau)$ and, by reality,
$\Phi_F(\tfrac12-i\tau)=\overline{\Phi_F(\tfrac12+i\tau)}$; hence
$g(\tfrac12+i\tau)=|\hat F(-\tau)|^2$, and $|\hat F|$ is even because $F$ is
real.  Estimate (8.2) is (8.1) applied to both factors, using
$|1-s-\tfrac12|=|s-\tfrac12|$.

For $h$: by definition
$h(a)=(F\star F^\dagger)(a)=\int F(t)\overline{F(t-a)}\,dt=\langle S_aF,F\rangle$,
which is real for real $F$ and satisfies $h(-a)=h(a)$; the support statement
and $\hat h=|\hat F|^2$ are standard properties of the convolution.  Finally
\[
  \Phi_{F^\dagger}(s)=\int\overline{F(-t)}e^{(s-\frac12)t}dt
  =\overline{\int F(u)e^{(\bar s-\frac12)(-u)}du}
  =\overline{\Phi_F(1-\bar s)}=\Phi_F(1-s)
\]
for real $F$, so $\Phi_h(s)=\Phi_F(s)\Phi_{F^\dagger}(s)=g(s)$.
\end{proof}

\subsection{The contour}

\begin{lemma}[the zero sum as a contour integral]
\label{lem:contour}
Fix $c>1$ and let $(X_j)_{j\ge1}$ be heights as in fact \emph{(E3b)}, i.e.\
$X_j\in[j,j+1]$ at distance at least $c_0/\log j$ from every ordinate of a zero
of $\xi$.  Let $R_j$ be the rectangle with vertices $1-c\mp iX_j$, $c\mp iX_j$,
traversed counterclockwise.  Then
\begin{equation}
  \sum_{|\Iml\rho|<X_j}g(\rho)
   =\frac1{2\pi i}\oint_{R_j}g(s)\frac{\xi'}{\xi}(s)\,ds
   \;\xrightarrow[j\to\infty]{}\;
   \frac1{2\pi i}\Bigl[\int_{L_c}-\int_{L_{1-c}}\Bigr]
     g(s)\frac{\xi'}{\xi}(s)\,ds ,
\tag{8.3}
\end{equation}
where $L_a$ is the line $\Rel s=a$ oriented upwards, both line integrals are
absolutely convergent, and the left-hand side converges to $\sum_\rho g(\rho)$
in the symmetric order $|\Iml\rho|\le X$, the zeros counted with multiplicity.
\end{lemma}

\begin{proof}
All zeros of $\xi$ lie in the strip $0\le\Rel s\le1$, which is contained in
$(1-c,c)$ because $c>1$; this is (E2) together with the nonvanishing of the
Euler product for $\Rel s>1$ and the functional equation.  The function $g$ is
entire (Lemma~\ref{lem:testfn}), so the only poles of $g\,\xi'/\xi$ inside
$R_j$ are the zeros of $\xi$ lying there, each a simple pole of $\xi'/\xi$ with
residue equal to its multiplicity.  The argument principle therefore gives the
first equality; no zero lies on $\partial R_j$, since $X_j$ avoids all
ordinates by (E3b) and the vertical sides $\Rel s=1-c,\,c$ lie outside the
critical strip.

\emph{Horizontal sides.}  By (E3c), $|\xi'/\xi(\sigma\pm iX_j)|=O(\log^2X_j)$
uniformly for $\sigma\in[1-c,c]$.  By (8.2) with $N=2$,
\[
  |g(\sigma\pm iX_j)|\le C_2^2e^{2T(c-\frac12)}X_j^{-4}
  \qquad(1-c\le\sigma\le c),
\]
since $|s-\tfrac12|\ge X_j$ there.  Each horizontal integral is therefore
$O\bigl(X_j^{-4}\log^2X_j\bigr)\to0$.

\emph{Vertical sides.}  On $L_c$ we have, from the logarithmic derivative of
(E2),
\[
  \frac{\xi'}{\xi}(s)=\frac1s+\frac1{s-1}-\tfrac12\log\pi
  +\tfrac12\psi(s/2)+\frac{\zeta'}{\zeta}(s),
\]
in which $\zeta'/\zeta$ is bounded on $L_c$ by (E1) and $\psi(s/2)=O(\log|s|)$
by (E4); hence $\xi'/\xi=O(\log|s|)$ on $L_c$, while $g$ decays faster than any
power of $|s|$ by (8.2).  The integral over $L_c$ is therefore absolutely
convergent, and the same holds on $L_{1-c}$ by the reflection
$\xi'/\xi(1-s)=-\xi'/\xi(s)$ of (E2).

Finally, the left-hand side of (8.3) is by construction the partial sum of the
zero series over $|\Iml\rho|<X_j$, that is, the symmetric partial sum.
\end{proof}

\begin{lemma}[the reflection]
\label{lem:reflection}
Let $A$ be meromorphic and integrable against $g$ on $L_c$ and $L_{1-c}$.  Then
\begin{equation}
  \Bigl[\int_{L_c}-\int_{L_{1-c}}\Bigr]g(s)A(s)\,ds
  =\int_{L_c}g(s)\bigl[A(s)-A(1-s)\bigr]\,ds .
\tag{8.4}
\end{equation}
In particular, since $\xi'/\xi(1-s)=-\xi'/\xi(s)$,
\begin{equation}
  \Bigl[\int_{L_c}-\int_{L_{1-c}}\Bigr]g\,\frac{\xi'}{\xi}\,ds
  =2\int_{L_c}g(s)\frac{\xi'}{\xi}(s)\,ds .
\tag{8.5}
\end{equation}
\end{lemma}

\begin{proof}
Substitute $s=1-u$ in the second integral.  As $s$ runs upwards along
$L_{1-c}$, the variable $u$ runs downwards along $L_c$, and $ds=-du$; the two
sign changes cancel, so
\[
  \int_{L_{1-c}}g(s)A(s)\,ds=\int_{L_c}g(1-u)A(1-u)\,du
  =\int_{L_c}g(u)A(1-u)\,du ,
\]
using $g(1-u)=g(u)$ (Lemma~\ref{lem:testfn}).  This is (8.4), and (8.5) is the
case $A=\xi'/\xi$.
\end{proof}

\subsection{Proof of the explicit formula}

\begin{theorem}[Theorem B]
\label{thm:EF}
Let $F\in C_c^\infty(\R)$ be real with $\supp F\subseteq[-T,T]$, put
$h=F\star F^\dagger$, and let the zero sum be taken in the symmetric order
$|\Iml\rho|\le X$.  Then
\begin{equation}
\begin{aligned}
  \sum_\rho\Phi_F(\rho)\Phi_F(1-\rho)
  ={}&2\Phi_F(0)\Phi_F(1)\\
  &+\frac1{2\pi}\int_\R|\hat F(\tau)|^2
     \bigl[\Rel\psi(\tfrac14+\tfrac{i\tau}2)-\log\pi\bigr]\,d\tau\\
  &-\sum_{2\le n<e^{2T}}\Lambda(n)n^{-1/2}
     \bigl[h(\log n)+h(-\log n)\bigr],
\end{aligned}
\tag{8.6}
\end{equation}
the prime sum being finite.
\end{theorem}

\begin{proof}
Write $g(s)=\Phi_F(s)\Phi_F(1-s)$ as in Lemma~\ref{lem:testfn}.  By
Lemma~\ref{lem:contour} and (8.5),
\begin{equation}
  \sum_\rho g(\rho)=\frac1{\pi i}\int_{L_c}g(s)\frac{\xi'}{\xi}(s)\,ds .
\tag{8.7}
\end{equation}
The reflection (8.5) must be applied before the logarithmic derivative is
expanded; see Remark~\ref{rem:route}(2).

On $L_c$ with $c>1$, taking the logarithmic derivative of
$\xi(s)=\tfrac12s(s-1)\pi^{-s/2}\Gamma(s/2)\zeta(s)$ gives
\begin{equation}
  \frac{\xi'}{\xi}(s)=A_0(s)+A_\Gamma(s)+\frac{\zeta'}{\zeta}(s),
  \qquad
  A_0(s):=\frac1s+\frac1{s-1},\quad
  A_\Gamma(s):=-\tfrac12\log\pi+\tfrac12\psi(s/2),
\tag{8.8}
\end{equation}
an identity of meromorphic functions, valid at every point where no term has a
pole, in particular on $L_c$.  Each of the three terms is integrable against
$g$ on $L_c$: $A_0$ is bounded there, $A_\Gamma=O(\log|s|)$ by (E4),
$\zeta'/\zeta$ is bounded by (E1), and $g$ decays faster than any power by
(8.2).  Hence (8.7) splits into three integrals, which we evaluate in turn.

\medskip
\noindent\textbf{(a) The polar term.}  Since
$A_0(1-s)=\frac1{1-s}+\frac1{-s}=-A_0(s)$, Lemma~\ref{lem:reflection} gives
\[
  \frac1{\pi i}\int_{L_c}gA_0\,ds
  =\frac1{2\pi i}\Bigl[\int_{L_c}-\int_{L_{1-c}}\Bigr]gA_0\,ds .
\]
The function $gA_0$ is meromorphic in the closed strip $1-c\le\Rel s\le c$
with exactly two poles, both simple, at $s=0$ and $s=1$ (note $1-c<0$ because
$c>1$), with residues $g(0)$ and $g(1)$ respectively.  Applying the residue
theorem on the rectangle $R_j$ of Lemma~\ref{lem:contour}, whose horizontal
contributions vanish in the limit by (8.2) and the boundedness of $A_0$ off the
poles, the bracket equals $2\pi i\,(g(0)+g(1))$.  Hence, by
Lemma~\ref{lem:testfn},
\begin{equation}
  \frac1{\pi i}\int_{L_c}gA_0\,ds=g(0)+g(1)=2\Phi_F(0)\Phi_F(1).
\tag{8.9}
\end{equation}

\medskip
\noindent\textbf{(b) The Gamma term.}  The function $A_\Gamma$ is holomorphic
in $\Rel s>0$ and, by (E4), $A_\Gamma(s)=O(\log|s|)$ there away from $s=0$.
Together with (8.2) this permits shifting the line $L_c$ to $L_{1/2}$: apply
Cauchy's theorem on the rectangle $[\tfrac12,c]\times[-X_j,X_j]$, whose
interior contains no pole of $gA_\Gamma$ and whose horizontal sides contribute
$O(X_j^{-4}\log X_j)\to0$ by (8.2) with $N=2$.  Hence, with $s=\tfrac12+i\tau$,
$ds=i\,d\tau$ and $\frac1{\pi i}(i\,d\tau)=\frac{d\tau}\pi$,
\[
  \frac1{\pi i}\int_{L_c}gA_\Gamma\,ds
  =\frac1{\pi}\int_\R g(\tfrac12+i\tau)
     \bigl[-\tfrac12\log\pi+\tfrac12\psi(\tfrac14+\tfrac{i\tau}2)\bigr]d\tau .
\]
Now $g(\tfrac12+i\tau)=|\hat F(\tau)|^2$ is even in $\tau$
(Lemma~\ref{lem:testfn}) and
$\psi(\tfrac14-\tfrac{i\tau}2)=\overline{\psi(\tfrac14+\tfrac{i\tau}2)}$, so
the odd imaginary part of the bracket integrates to zero and
\begin{equation}
  \frac1{\pi i}\int_{L_c}gA_\Gamma\,ds
  =\frac1{2\pi}\int_\R|\hat F(\tau)|^2
     \bigl[\Rel\psi(\tfrac14+\tfrac{i\tau}2)-\log\pi\bigr]d\tau .
\tag{8.10}
\end{equation}

\medskip
\noindent\textbf{(c) The prime terms.}  On $L_c$ fact (E1) gives
$\zeta'/\zeta(s)=-\sum_{n\ge2}\Lambda(n)n^{-s}$ with
$\sum_n\Lambda(n)n^{-c}<\infty$; since $\int_{L_c}|g|<\infty$ by
Lemma~\ref{lem:contour}, Fubini's theorem permits term-by-term integration:
\[
  \frac1{\pi i}\int_{L_c}g\frac{\zeta'}{\zeta}\,ds
  =-\sum_{n\ge2}\Lambda(n)\cdot\frac1{\pi i}\int_{L_c}g(s)n^{-s}\,ds .
\]
The integrand $g(s)n^{-s}$ is now \emph{entire}: $g$ is entire by
Lemma~\ref{lem:testfn} and $n^{-s}=e^{-s\log n}$ is entire.  The logarithmic
derivative $\zeta'/\zeta$ has already been expanded and is no longer present,
so the line may be shifted from $L_c$ to $L_{1/2}$ with no residue of any kind
and, in particular, without crossing a zero of $\zeta$.  (Apply Cauchy's
theorem on $[\tfrac12,c]\times[-X_j,X_j]$; the horizontal contributions are
$O(X_j^{-4})\to0$ by (8.2), uniformly in $n\ge2$ because
$|n^{-s}|=n^{-\sigma}\le n^{-1/2}$ on the strip.)  Hence, with
$s=\tfrac12+i\tau$ and $g(\tfrac12+i\tau)=\hat h(\tau)$,
\[
  \frac1{\pi i}\int_{L_c}g(s)n^{-s}\,ds
  =2n^{-1/2}\cdot\frac1{2\pi}\int_\R\hat h(\tau)e^{-i\tau\log n}\,d\tau
  =2n^{-1/2}h(-\log n),
\]
by Fourier inversion (E5).  Since $h$ is even (Lemma~\ref{lem:testfn}, where
this uses that $F$ is real), $2h(-\log n)=h(\log n)+h(-\log n)$; and the sum is
finite because $\supp h\subseteq[-2T,2T]$, so that only $n<e^{2T}$ contributes.
Therefore
\begin{equation}
  \frac1{\pi i}\int_{L_c}g\frac{\zeta'}{\zeta}\,ds
  =-\sum_{2\le n<e^{2T}}\Lambda(n)n^{-1/2}\bigl[h(\log n)+h(-\log n)\bigr].
\tag{8.11}
\end{equation}

Adding (8.9), (8.10) and (8.11) and recalling
$\sum_\rho g(\rho)=\sum_\rho\Phi_F(\rho)\Phi_F(1-\rho)$ gives (8.6).
\end{proof}

\begin{corollary}
\label{cor:EFeven}
Let $h\in C_c^\infty(\R)$ be real and even with $\supp h\subseteq[-2T,2T]$, and
put $g:=\Phi_h$.  Then $g$ is entire, $g(1-s)=g(s)$, $g(\tfrac12+i\tau)=\hat
h(\tau)$ is real (not necessarily nonnegative), and
\begin{equation}
\begin{aligned}
  \sum_\rho\Phi_h(\rho)={}&\Phi_h(0)+\Phi_h(1)
   +\frac1{2\pi}\int_\R\hat h(\tau)
     \bigl[\Rel\psi(\tfrac14+\tfrac{i\tau}2)-\log\pi\bigr]d\tau\\
   &-\sum_{2\le n<e^{2T}}\Lambda(n)n^{-1/2}\bigl[h(\log n)+h(-\log n)\bigr].
\end{aligned}
\tag{8.12}
\end{equation}
Theorem~\emph{\ref{thm:EF}} is the case $h=F\star F^\dagger$, in which
$g(s)=\Phi_F(s)\Phi_F(1-s)$ and $\hat h=|\hat F|^2$.
\end{corollary}

\begin{proof}
The proof of Theorem~\ref{thm:EF} is repeated with
Lemma~\ref{lem:testfn} replaced by the following elementary facts.  $\Phi_h$ is
entire and satisfies (8.1) with $C_N=\|h^{(N)}\|_{L^1}$ and $T$ replaced by
$2T$, by the same integration by parts.  The symmetry $g(1-s)=g(s)$ \emph{is}
the evenness of $h$:
\[
  \Phi_h(1-s)=\int h(t)e^{(\frac12-s)t}dt=\int h(-u)e^{(s-\frac12)u}du
  =\Phi_h(s).
\]
And $g(\tfrac12+i\tau)=\int h(t)e^{i\tau t}dt=\hat h(-\tau)=\hat h(\tau)$,
which is real because $h$ is real and even.  Step (a) of the proof uses only
$g(0)=\Phi_h(0)$ and $g(1)=\Phi_h(1)$; step (b) uses only that
$g(\tfrac12+i\cdot)$ is real and even, which is where the positivity of
$|\hat F|^2$ was never needed; step (c) uses only Fourier inversion and the
evenness of $h$.  Lemmas~\ref{lem:contour} and \ref{lem:reflection} are
unchanged, since they use only that $g$ is entire, that $g(1-s)=g(s)$, and that
$|g(\sigma+i\tau)|\le C_4\,e^{2T|\sigma-\frac12|}\,|s-\tfrac12|^{-4}$, which is
(8.1) with $N=4$; that bound is all that the horizontal estimate of
Lemma~\ref{lem:contour} uses, and it replaces (8.2), which was derived from a
product of two transforms and is not available for a single one.
\end{proof}

\begin{remark}[absolute convergence]
\label{rem:abs}
For $h\in C_c^\infty(\R)$ real and even, (8.1) applied to $\Phi_h$ with $N=2$
gives $|\Phi_h(\rho)|\le C_2e^{2T|\beta-\frac12|}|\rho-\tfrac12|^{-2}
=O(|\gamma_\rho|^{-2})$ uniformly for $\rho=\beta+i\gamma_\rho$ with
$0\le\beta\le1$ and $|\gamma_\rho|\ge1$; and
$\sum_{|\gamma_\rho|\ge1}|\gamma_\rho|^{-2}<\infty$ by partial summation
against $N(X)=O(X\log X)$, fact (E3d).  The zeros with $|\gamma_\rho|<1$ are
finite in number, again by (E3d), and $\Phi_h$ is entire, so they contribute a
finite amount.  Hence $\sum_\rho\Phi_h(\rho)$ converges absolutely.  The
symmetric order is needed for the
argument-principle step of Lemma~\ref{lem:contour}, and is stated there; it
imposes no restriction on the \emph{value} of the sum.
\end{remark}

\begin{remark}[on the proof]
\label{rem:route}
\begin{enumerate}
\item[(1)] The single integral $\int_{L_c}g(s)n^{-s}ds$, shifted to
$\Rel s=\tfrac12$, produces $n^{-1/2}h(-\log n)$, which is \emph{half} of the
symmetric prime term.  The other half comes from the line $L_{1-c}$ through
(8.5), and the two halves coincide only because $h$ is even, that is, only
because $F$ is real.  Both facts are used explicitly above.
\item[(2)] The termwise split (8.8) must be made on $L_c$ \emph{only}.  One may
not split the difference $\bigl[\int_{L_c}-\int_{L_{1-c}}\bigr]$ termwise,
because the constant $-\tfrac12\log\pi$ is entire and symmetric, so that
combination would annihilate exactly the $\log\pi$ the formula needs.  This is
why (8.5) is applied first and the expansion second.
\item[(3)] The horizontal estimate requires $N\ge2$ in (8.1), i.e.\ two
derivatives of $F$; $F\in C_c^\infty$ is ample, but a merely continuous
compactly supported $F$ would not do.
\item[(4)] No contour is shifted across a zero anywhere in the proof.  In
step (c) the integrand is entire, so the shift is free; in step (b) the
integrand is holomorphic in $\Rel s\ge\tfrac12$; and in step (a) the residues
crossed are exactly the two poles of the trivial factor $\tfrac12s(s-1)$ of
$\xi$, computed explicitly.
\end{enumerate}
\end{remark}

\section{The comparison theorem}
\label{sec:comparison}

\subsection{The Weil quadratic form}

\begin{definition}
\label{def:weilform}
For real $F\in C_c^\infty(\R)$ with $h=F\star F^\dagger$ put
\begin{equation}
  \mathcal W(F):=\frac1{2\pi}\int_\R|\hat F(\tau)|^2
     \bigl[\Rel\psi(\tfrac14+\tfrac{i\tau}2)-\log\pi\bigr]d\tau
   -\sum_{n\ge2}\Lambda(n)n^{-1/2}\bigl[h(\log n)+h(-\log n)\bigr],
\tag{9.1}
\end{equation}
the prime sum being finite.  Equivalently
$\mathcal W(F)=\langle\mathcal AF,F\rangle$ for the operator
\[
  \mathcal A=g_\Gamma(D)-m_0-\sum_{n\ge2}\Lambda(n)n^{-1/2}
   \bigl(S_{\log n}+S_{-\log n}\bigr),
  \qquad g_\Gamma(\tau)-m_0=\Rel\psi(\tfrac14+\tfrac{i\tau}2)-\log\pi ,
\]
with $(S_aF)(t)=F(t+a)$ and $g_\Gamma(D)$ the Fourier multiplier by
$g_\Gamma$; this is the classical Weil quadratic form of the test $F$.  In this
notation Theorem~\emph{\ref{thm:EF}} reads
\begin{equation}
  \sum_\rho\Phi_F(\rho)\Phi_F(1-\rho)=2\Phi_F(0)\Phi_F(1)+\mathcal W(F).
\tag{9.2}
\end{equation}
\end{definition}

\begin{lemma}
\label{lem:traceweil}
Let $F\in C_c^\infty(\R)$ be real, $h=F\star F^\dagger$, and let $T_\infty$ be
the colimit trace of Proposition~\emph{\ref{prop:colimit}} with $c=0$.  Then
\begin{equation}
  I_\infty(\Gamma_F,\Gamma_F)=T_\infty(\Gamma_h)=-\mathcal W(F).
\tag{9.3}
\end{equation}
At a fixed finite chart $S$ the same identity holds provided
$R(\Gamma_h)<\log\ell_{\mathrm{next}}(S)$.
\end{lemma}

\begin{proof}
By Proposition~\ref{prop:radalg},
$\Gamma_F\star\Gamma_F^\dagger=\Gamma_{F\star F^\dagger}=\Gamma_h$, so
$I_\infty(\Gamma_F,\Gamma_F)=T_\infty(\Gamma_h)$.  Note that $h(0)=\|F\|_{L^2}^2>0$,
so $0\in\supp h$ and Proposition~\ref{prop:primepacket} is not applicable; we
use Lemma~\ref{lem:radial-trace}, which carries no support hypothesis, in its
colimit form.  By Lemma~\ref{lem:testfn} $h$ is real and even, so
$\hat h(-\tau)=\hat h(\tau)=|\hat F(\tau)|^2$; and
$\widehat W_\infty(0,\tau)=\log\pi-\Rel\psi(\tfrac14+\tfrac{i\tau}2)$ by
Theorem~\ref{thm:archsymbol}.  Hence (4.3) reads
\[
  T_\infty(\Gamma_h)=\sum_{n\ge2}\Lambda(n)n^{-1/2}
     \bigl[h(\log n)+h(-\log n)\bigr]
   +\frac1{2\pi}\int_\R|\hat F(\tau)|^2
     \bigl[\log\pi-\Rel\psi(\tfrac14+\tfrac{i\tau}2)\bigr]d\tau ,
\]
both terms absolutely convergent, the first a finite sum, where we have
rewritten the double sum over $p$ and $k\ge1$ as the sum over prime powers
$n=p^k$ using $\Lambda(p^k)=\log p$.  This is exactly $-\mathcal W(F)$.  The
chart clause is Proposition~\ref{prop:colimit}.
\end{proof}

\subsection{The comparison theorem}

\begin{theorem}[Theorem C]
\label{thm:comparison}
Let $F\in C_c^\infty(\R)$ be real, $h=F\star F^\dagger$, and let
$\Gamma_F,\Gamma_h\in\DAr$ be the corresponding radial classes.  Let $c=0$ and
let $T_\infty$ be the colimit trace.  Then, the zero sum in symmetric order,
\begin{equation}
\begin{aligned}
  I_\infty(\Gamma_F,\Gamma_F)=T_\infty(\Gamma_{F\star F^\dagger})
  &=2\Phi_F(0)\Phi_F(1)-\sum_\rho\Phi_F(\rho)\Phi_F(1-\rho)\\
  &=\beta_T\bigl(M(\Gamma_F),M(\Gamma_F)\bigr)
     -\sum_\rho\Phi_F(\rho)\Phi_F(1-\rho),
\end{aligned}
\tag{9.4}
\end{equation}
and consequently, when $\Gamma_F$ is primitive, that is, when $M(\Gamma_F)=0$,
\begin{equation}
  I_\infty(\Gamma_F,\Gamma_F)=-\sum_\rho\Phi_F(\rho)\Phi_F(1-\rho)
  =-\mathcal W(F).
\tag{9.5}
\end{equation}
At a fixed finite chart $S$ identities \emph{(9.4)}--\emph{(9.5)} hold
provided $R(\Gamma_h)<\log\ell_{\mathrm{next}}(S)$; outside that window the
left-hand side omits the contributions of the primes $\ell\notin S_f$ and the
identity fails.
\end{theorem}

\begin{proof}
Combine (9.2) with (9.3): $\mathcal W(F)=\sum_\rho\Phi_F(\rho)\Phi_F(1-\rho)
-2\Phi_F(0)\Phi_F(1)$, so
$T_\infty(\Gamma_h)=-\mathcal W(F)=2\Phi_F(0)\Phi_F(1)
-\sum_\rho\Phi_F(\rho)\Phi_F(1-\rho)$.  The identification of the polar term
with $\beta_T(M(\Gamma_F),M(\Gamma_F))$ is Theorem~\ref{thm:polarpacket}(i),
and (9.5) is the case $M(\Gamma_F)=0$, which by
Theorem~\ref{thm:polarpacket}(ii) is primitivity.  The chart clause is
Proposition~\ref{prop:colimit} and the chart clause of
Lemma~\ref{lem:traceweil}.
\end{proof}

\begin{corollary}[Theorem D]
\label{cor:c-unique}
Let $c\in\R$ and let $T^{(c)}=\sum_vW_v+c\,(\,\cdot\,)(1)$ be the corresponding
member of the family of Proposition~\emph{\ref{prop:trace}}, with
$I^{(c)}(F,G)=T^{(c)}(F\star G^\dagger)$ and $I_\infty^{(c)}$ its colimit.
Then for every real $F\in C_c^\infty(\R)$:
\begin{enumerate}
\item[(i)] $\displaystyle I_\infty^{(c)}(\Gamma_F,\Gamma_F)
   =2\Phi_F(0)\Phi_F(1)-\sum_\rho\Phi_F(\rho)\Phi_F(1-\rho)
   +c\,\|F\|_{L^2(\R)}^2$;
\item[(ii)] for $F\ne0$ the correction $c\|F\|_{L^2}^2$ is strictly positive if
$c>0$ and strictly negative if $c<0$, and it does \emph{not} vanish on the
primitive classes: the condition $M(\Gamma_F)=0$ annihilates the polar term,
not the $L^2$-norm;
\item[(iii)] consequently $c=0$ is the unique real constant for which
\emph{(9.4)} holds for all real $F\in C_c^\infty(\R)$.
\end{enumerate}
\end{corollary}

\begin{proof}
(i)  By definition $T^{(c)}=T+c\,(\,\cdot\,)(1)$, so
$I^{(c)}(\Gamma_F,\Gamma_F)=I_A(\Gamma_F,\Gamma_F)
+c\,(\Gamma_F\star\Gamma_F^\dagger)(1)$, and
\[
  (\Gamma_F\star\Gamma_F^\dagger)(1)
  =\int_{\CS}\Gamma_F(u)\overline{\Gamma_F(u)}\,d^*u=\|\Gamma_F\|_{L^2(\CS)}^2 .
\]
By Lemma~\ref{lem:profile} and the disintegration $d^*w=dt\otimes\mu$ with
$\mu$ a probability measure,
$\|\Gamma_F\|_{L^2(\CS)}^2=\int_\R|F(-t)|^2dt=\|F\|_{L^2(\R)}^2$.  Adding this
term to (9.4) gives the displayed identity.  (Equivalently, on the principal
component the Plancherel form (2.6) contributes
$\frac c{2\pi}\int|\hat F|^2d\tau=c\|F\|_{L^2}^2$ by (E5).)

(ii)  $\|F\|_{L^2}^2>0$ for $F\ne0$, and nothing in
$M(\Gamma_F)=(\Phi_F(1),\Phi_F(0))$ constrains it: for instance the primitive
tests $F=\Pi G$ constructed in Corollary~\ref{cor:Pi} below are nonzero and
primitive, so they have $\|F\|_{L^2}^2>0$ and $M(\Gamma_F)=0$ simultaneously.

(iii)  Two values $c\ne c'$ give values of $I_\infty^{(c)}(\Gamma_F,\Gamma_F)$
differing by $(c-c')\|F\|_{L^2}^2\ne0$ for $F\ne0$, so at most one $c$ can
satisfy (9.4); and $c=0$ does, by Theorem~\ref{thm:comparison}.
\end{proof}

\begin{remark}
\label{rem:loadbearing}
The identities of Sections~\ref{sec:comparison}--\ref{sec:sufficiency} concern
the trace with $c=0$: for any other $c$ the right-hand side of (9.4) acquires
the term $c\|F\|_{L^2}^2$, which does not vanish on primitive classes.  We do
not consider here whether an inequality $I_\infty^{(c)}\le0$ on primitive
radial classes would, for $c\ne0$, still be equivalent to a statement about the
zeros.
\end{remark}

\subsection{The limit of the chart symbols}

At a finite chart the symbol $\widehat T_S(\cz,\cdot)$ is a
real-analytic function (Proposition~\ref{prop:symbol}).  In the colimit there
is no such function; what survives is a linear functional, and the explicit
formula says exactly what it is.

\begin{theorem}
\label{thm:limitsymbols}
Let $\mathcal E:=C_c^\infty(\R)^{\mathrm{even},\,\mathrm{real}}$, a real vector
space, and for a chart $S$ and $h\in\mathcal E$ put
\[
  L_S(h):=\frac1{2\pi}\int_\R\hat h(\tau)\,\widehat T_S(\cz,\tau)\,d\tau
  =T_S(\Gamma_h).
\]
Then, as $S$ increases along any exhaustion of the primes, the sequence
$L_S(h)$ is \emph{eventually constant} for every $h\in\mathcal E$ --- constant
for every $S$ containing all primes $\ell\le e^{R(\Gamma_h)}$ --- with common
value $L_\infty(h)=T_\infty(\Gamma_h)$, and
\begin{equation}
  L_\infty(h)=\Phi_h(0)+\Phi_h(1)-\sum_\rho\Phi_h(\rho),
  \qquad h\in\mathcal E ,
\tag{9.6}
\end{equation}
a linear functional on $\mathcal E$.  Theorem~\emph{\ref{thm:comparison}} is
the case $h=F\star F^\dagger$.
\end{theorem}

\begin{proof}
$L_S(h)=T_S(\Gamma_h)$ is Lemma~\ref{lem:radial-trace} together with the
Plancherel form of Proposition~\ref{prop:symbol}(c) and
Lemma~\ref{lem:profile}, using that $\hat h$ is even and real for
$h\in\mathcal E$.  Linearity of $h\mapsto L_\infty(h)$ is linearity of
$T_\infty$ and of $h\mapsto\Gamma_h$.  Eventual constancy is
Proposition~\ref{prop:colimit}: $W_\ell(J_{ST}\Gamma_h)=0$ for every prime
$\ell$ with $\log\ell>R(\Gamma_h)$, so the sequence is constant for every $S$
containing the finitely many primes $\ell\le e^{R(\Gamma_h)}$, with common
value $T_\infty(\Gamma_h)$; the convergence is eventually exact, not merely
convergent.  The value (9.6) is Corollary~\ref{cor:EFeven} rearranged: add
$\Phi_h(0)+\Phi_h(1)$ to (8.12) and use Lemma~\ref{lem:radial-trace}.
\end{proof}

\begin{remark}
\label{rem:notafunction}
The convergence in Theorem~\ref{thm:limitsymbols} is eventual constancy for
each fixed $h$; no continuity of $L_\infty$ on $\mathcal E$ is asserted.  The
limit functional is not given by integration against a function: at $\tau=0$ the finite-place part of the colimit symbol
diverges,
\[
  \sum_p(\log p)\,\frac{2p^{-1/2}}{1-p^{-1/2}}=+\infty
\]
(\cite{P44}, Remark 9.12), so there is no function on the principal
component against which $L_\infty$ integrates.
\end{remark}

\begin{remark}
\label{rem:noposition}
No step of Sections~\ref{sec:EF}--\ref{sec:comparison} uses the location of
the zeros.  Fact (E2) is used only through the statement that the zeros of $\xi$ are
the nontrivial zeros of $\zeta$ and lie in the critical strip; (E3a)--(E3d) use
only their density and the Hadamard product; Lemma~\ref{lem:contour} chooses
heights $X_j$ avoiding the ordinates, which requires no knowledge of where
those ordinates are; and the contour shifts of Theorem~\ref{thm:EF} are
performed on integrands that are entire or holomorphic in the relevant
half-plane.
\end{remark}

\section{Residues of the Euler-product symbol}
\label{sec:polar}

The polar term $2\Phi_F(0)\Phi_F(1)$ of Theorem~\ref{thm:comparison} arose in
step (a) of the proof of Theorem~\ref{thm:EF} from two residues.  In this
section we identify them as the residues of a single meromorphic function and
compare them with the residues of the chart symbols.  At $s=0$ the chart
residues grow with $S$ and the surviving residue is the archimedean one; at
$s=1$ no chart symbol has a pole, and the pole appears only in the limit.

\subsection{The local symbols in the \texorpdfstring{$s$}{s}-variable}

\begin{lemma}
\label{lem:localsymbols}
Put $s=\tfrac12+i\tau$ and, for a prime $p$ and for the real place,
\begin{equation}
  f_p(s):=\frac{(\log p)\,p^{-s}}{1-p^{-s}}=\sum_{k\ge1}(\log p)p^{-ks}
  \quad(\Rel s>0),\qquad
  a(s):=\tfrac12\log\pi-\tfrac12\psi(s/2).
\tag{10.1}
\end{equation}
Then, as meromorphic functions of $s$, on the principal component $\cz$:
\begin{enumerate}
\item[(i)] $\widehat W_p(\cz,\tau)=f_p(s)+f_p(1-s)$;
\item[(ii)] $\widehat W_\infty(0,\tau)=a(s)+a(1-s)$;
\item[(iii)] $\Res_{s=0}f_p=\Res_{s=0}a=1$, while
$\Res_{s=1}f_p(1-s)=\Res_{s=1}a(1-s)=-1$; moreover $f_p$ has no pole at $s=1$
and $a$ has no pole in $\Rel s>-2$ other than the one at $s=0$;
\item[(iv)] consequently, at the chart $S$,
$\widehat T_S(\cz,\cdot)=A_S(s)+A_S(1-s)$ with
$A_S(s):=a(s)+\sum_{p\in S_f}f_p(s)$, and
\[
  \Res_{s=0}A_S=|S_f|+1 .
\]
\end{enumerate}
\end{lemma}

\begin{proof}
(i)  By (2.5), on $\cz$ one has
$\widehat W_p=2(\log p)\Rel\bigl(\zeta_p/(1-\zeta_p)\bigr)$ with
$\zeta_p=p^{-1/2+i\tau}=p^{-(1-s)}$.  Hence
$(\log p)\zeta_p/(1-\zeta_p)=f_p(1-s)$, and for real $\tau$,
$\overline{f_p(1-s)}=f_p(\overline{1-s})=f_p(s)$, so that
$2\Rel(\cdot)=f_p(s)+f_p(1-s)$.  Both sides are meromorphic and agree on a
line, hence everywhere.

(ii) is Theorem~\ref{thm:archresidue}(i).

(iii)  From $1-p^{-s}=s\log p+O(s^2)$ one gets $f_p(s)=1/s+O(1)$, so
$\Res_{s=0}f_p=1$; the poles of $f_p$ are the points $s\in(2\pi i/\log p)\Z$,
so $s=1$ is not among them.  For $a$ this is
Theorem~\ref{thm:archresidue}(ii).  The residues of $s\mapsto f_p(1-s)$ and
$s\mapsto a(1-s)$ at $s=1$ are the negatives of the residues of $f_p,a$ at
$s=0$, by the chain rule.

(iv)  Sum (i) over $p\in S_f$ and add (ii).  At $s=0$ the residue receives $1$
from $a$ and $1$ from each $f_p$, $p\in S_f$, while $f_p(1-s)$ and $a(1-s)$ are
regular there, taking the finite values $f_p(1)=(\log p)/(p-1)$ and $a(1)$.
\end{proof}

\subsection{The poles at \texorpdfstring{$s=0$ and $s=1$}{s=0 and s=1}}

\begin{definition}
\label{def:A}
Put
\begin{equation}
  A(s):=a(s)+\sum_pf_p(s)
       =\tfrac12\log\pi-\tfrac12\psi(s/2)-\frac{\zeta'}{\zeta}(s),
\tag{10.2}
\end{equation}
the sum over all primes, absolutely convergent for $\Rel s>1$ by (E1) and
continued meromorphically to $\C$ by the right-hand expression; and, at a
finite chart, $A_S(s):=a(s)+\sum_{p\in S_f}f_p(s)$, so that
$\widehat T_S(\cz,\cdot)=A_S(s)+A_S(1-s)$ by
Lemma~\ref{lem:localsymbols}(iv).
\end{definition}

Note that $A$ is not the symbol of the colimit trace $T_\infty$.  By
Theorem~\ref{thm:limitsymbols} and Remark~\ref{rem:notafunction}, $T_\infty$ is
a finite sum at each test and has no Fourier-side function at all; $A$ is the
analytic continuation of a prime sum that does not converge on the critical
line, where the chart symbols live; compare
Corollary~\ref{cor:noncommute} below.

\begin{theorem}[Theorem E]
\label{thm:residues}
With the notation above:
\begin{enumerate}
\item[(i)] $\displaystyle A(s)=-\frac{\xi'}{\xi}(s)+\frac1s+\frac1{s-1}$;
\item[(ii)] $\xi$ is entire and nonzero at $s=0$ and $s=1$, so $-\xi'/\xi$ is
regular there and $A$ has exactly two poles in $\{0,1\}$, both simple, with
\begin{equation}
  \Res_{s=0}A=1,\qquad \Res_{s=1}A=1 ,
\tag{10.3}
\end{equation}
chart-independently; they are the poles of the trivial factor $\tfrac12s(s-1)$
of $\xi$ and of nothing else;
\item[(iii)] at every finite chart, by contrast,
\begin{equation}
  \Res_{s=0}A_S=|S_f|+1,\qquad \Res_{s=1}A_S=0 ;
\tag{10.4}
\end{equation}
\item[(iv)] in step \emph{(a)} of the proof of Theorem~\emph{\ref{thm:EF}} the
two summands of $A_0(s)=1/s+1/(s-1)$ contribute, against $g$, exactly the
residues $g(0)$ and $g(1)$; by \emph{(i)} these are the same two poles as in
\emph{(ii)}, with the same residues $1$ each.  Hence
\[
  g(0)+g(1)=2\Phi_F(0)\Phi_F(1)
  =\beta_T\bigl(M(\Gamma_F),M(\Gamma_F)\bigr),
\]
that is, the polar term is given by the two poles of $A$, with coefficient
exactly $1$ at each, the pole at $s=0$ corresponding under
Theorem~\emph{\ref{thm:polarpacket}} and Remark~\emph{\ref{rem:swap}} to the
isotropic direction $e_+$ and the pole at $s=1$ to $e_-$.
\end{enumerate}
\end{theorem}

\begin{proof}
(i)  Logarithmic differentiation of
$\xi(s)=\tfrac12s(s-1)\pi^{-s/2}\Gamma(s/2)\zeta(s)$ gives (8.8),
\[
  \frac{\xi'}{\xi}(s)=\frac1s+\frac1{s-1}-\tfrac12\log\pi+\tfrac12\psi(s/2)
  +\frac{\zeta'}{\zeta}(s),
\]
whence
$\tfrac12\log\pi-\tfrac12\psi(s/2)-\zeta'/\zeta(s)
=-\xi'/\xi(s)+1/s+1/(s-1)$, which is (i) by (10.2).

(ii)  By (E2) the function $\xi$ is entire and $\xi(0)=\xi(1)=\tfrac12\ne0$:
the zero of $s-1$ cancels the pole of $\zeta$ at $s=1$, the zero of $s$ cancels
the pole of $\Gamma(s/2)$ at $s=0$, and the functional equation then gives
$\xi(0)=\xi(1)$.  So $\xi'/\xi$ is regular at $0$ and $1$, and by (i) the only
poles of $A$ there are $1/s$ and $1/(s-1)$, each simple with residue $1$.  The
same two residues may be read off (10.2) directly, which is a useful
cross-check: at $s=0$ the function $\zeta'/\zeta$ is regular, since
$\zeta(0)=-\tfrac12\ne0$, and $-\tfrac12\psi(s/2)=1/s+O(1)$ by (E4), so the
residue is the \emph{archimedean} one; at $s=1$ the value $\psi(\tfrac12)$ is
finite while $-\zeta'/\zeta(s)=1/(s-1)+O(1)$ by the simple pole of $\zeta$, so
the residue comes from the finite places, that is, from the whole Euler
product at once.

(iii)  The first assertion is Lemma~\ref{lem:localsymbols}(iii),(iv).  For the
second, $a$ is regular at $s=1$ and each $f_p$ has poles only on
$(2\pi i/\log p)\Z$, so $f_p(1)=(\log p)/(p-1)<\infty$ and the finite sum
$A_S(1)=a(1)+\sum_{p\in S_f}(\log p)/(p-1)$ is finite: a finite chart has no
pole at $s=1$ at all.

(iv)  Step (a) of the proof of Theorem~\ref{thm:EF} computed
$\frac1{\pi i}\int_{L_c}gA_0\,ds=g(0)+g(1)$ by the residue theorem on $R_j$,
and $A_0$ is exactly the two-pole part isolated in (i)--(ii).  By
Lemma~\ref{lem:testfn}, $g(0)=g(1)=\Phi_F(0)\Phi_F(1)$, and
Theorem~\ref{thm:polarpacket}(i) identifies the sum with
$\beta_T(M(\Gamma_F),M(\Gamma_F))$.  The correspondence of the two poles with
the two isotropic directions is Remark~\ref{rem:swap}.
\end{proof}

\begin{remark}
\label{rem:collapses}
The comparison of (10.3) with (10.4) is asymmetric.
\begin{center}
\begin{tabularx}{\textwidth}{@{}lllY@{}}
\toprule
point & at a finite chart $S$ & in the continuation & what the passage is\\
\midrule
$s=0$ & $\Res=|S_f|+1$, divergent as $S\uparrow$ & $\Res=1$ &
the $|S_f|$ finite residues disappear; the surviving $1$ is the archimedean
one, and it is the $1/s$ of the trivial factor of $\xi$
(Corollary~\ref{cor:noncommute})\\[2pt]
$s=1$ & $\Res=0$ for every $S$ & $\Res=1$ &
the pole is created in the limit, by the pole of $\zeta$ at $s=1$;
the rate at which the chart values approach it is Mertens'
(Proposition~\ref{prop:mertens})\\
\bottomrule
\end{tabularx}
\end{center}
In both cases the continuation has residue exactly $1$, and the two residues
together give the polar term.
\end{remark}

\begin{remark}
\label{rem:reflected}
Write $\Xi(s):=A(s)+A(1-s)$.  Then $\Xi\equiv0$ on $\C$: by
Theorem~\ref{thm:residues}(i) and $\xi'/\xi(1-s)=-\xi'/\xi(s)$,
\[
  A(1-s)=\frac{\xi'}{\xi}(s)+\frac1{1-s}+\frac1{-s}
        =\frac{\xi'}{\xi}(s)-\frac1{s-1}-\frac1s=-A(s).
\]
In particular the residue $+1$ of $A$ at $s=0$ is cancelled inside $\Xi$ by the
residue $-1$ of $A(1-s)$ there, the latter coming from the pole of $\zeta$ at
$s=1$.  This identity is a restatement of the functional equation, and
nothing below uses it.
\end{remark}

\begin{corollary}
\label{cor:noncommute}
The map $S\mapsto\Res_{s=0}A_S=|S_f|+1$ diverges as $S\uparrow$, while
$\Res_{s=0}A=1$: summation over the primes and analytic continuation do not
commute at $s=0$.  Precisely,
$\sum_{p\in S_f}f_p(s)\to-\zeta'/\zeta(s)$ only on $\Rel s>1$; at $s=0$ each
$f_p$ has a pole of residue $1$ while $-\zeta'/\zeta$, the analytic
continuation of the sum, is regular there.  The omitted tail
$\sum_{p\notin S_f}f_p$ therefore supplies exactly one $1/s$ per omitted prime
near $s=0$, the sequence of chart residues is unbounded, and the continuation
contributes none of them.
\end{corollary}

\begin{proof}
Lemma~\ref{lem:localsymbols}(iv) and
Theorem~\ref{thm:residues}(ii),(iii); the divergence of the tail residues is
the count of omitted primes.
\end{proof}

\begin{proposition}
\label{prop:mertens}
For $S_f=\{p\le x\}$,
\[
  A_S(1)=a(1)+\sum_{p\le x}\frac{\log p}{p-1}=\log x+O(1),
\]
while $-\zeta'/\zeta(\sigma)=1/(\sigma-1)+O(1)$ as $\sigma\downarrow1$.  Thus
Mertens' theorem governs the behaviour of the chart symbols at $s=1$, that is,
at $\tau=-i/2$; it has nothing to do with the residue count at $s=0$, which is
the non-commutation of Corollary~\emph{\ref{cor:noncommute}}.
\end{proposition}

\begin{proof}
By fact (E8), $\sum_{N\le x}\Lambda(N)/N=\log x+O(1)$.  Since
\[
  \sum_{p\le x}\frac{\log p}{p-1}=\sum_{p\le x}\sum_{k\ge1}\frac{\log p}{p^k}
  =\sum_{N\le x}\frac{\Lambda(N)}N+O(1),
\]
the higher prime powers contributing $O(1)$ in total, the first assertion
follows, $a(1)$ being a fixed finite constant.  The second is the simple pole
of $\zeta$ at $s=1$, fact (E2).
\end{proof}

\begin{remark}
\label{rem:mertens-scope}
Fact (E8) is used only in Proposition~\ref{prop:mertens}, and through
\cite{P43}, Theorem 15.2, in the following statement, which we quote rather
than reprove: in the colimit the form $I_\infty$ restricted to the plane of
isotropic directions reproduces $\beta_T$, with structure constant
$\beta_T(e_+,e_-)=1$ arising from the main term $\log x$ of Mertens' theorem.
It shows that the hyperbolic form $\beta_T$ in (9.4) is the restriction of the
colimit form to that plane.  It is not used to identify the polar term, which
is Theorem~\ref{thm:residues}(iv), and none of the main theorems depends on
it.
\end{remark}

\section{Primitive test functions detect zeros off the line}
\label{sec:sufficiency}

Theorem~\ref{thm:comparison} identifies the trace on primitive radial classes
with $-\sum_\rho\Phi_F(\rho)\Phi_F(1-\rho)$.  The test functions in (9.5) are
subject to the two linear conditions $\Phi_F(0)=\Phi_F(1)=0$, and we show that
they still detect a zero off the critical line.  The reason is that the
conditions are imposed by multiplication by an explicit real even polynomial
which does not vanish at any zero of $\xi$.

\subsection{The zero coordinate}

Write
\begin{equation}
  \mathcal Z:=\{z\in\C:\ \xi(\tfrac12+iz)=0\},
\tag{11.1}
\end{equation}
with multiplicity; thus $\rho=\tfrac12+iz$ is a bijection between $\mathcal Z$
and the zeros of $\xi$.  Since all zeros lie in $0\le\Rel s\le1$ we have
$|\Iml z|\le\tfrac12$ on $\mathcal Z$, and the divisor $\mathcal Z$ is
invariant under $z\mapsto-z$ (from $\rho\mapsto1-\rho$) and under
$z\mapsto\bar z$ (from $\rho\mapsto\bar\rho$ composed with the former).

For $F\in C_c^\infty(\R)$ put
\begin{equation}
  \mathfrak F_F(z):=\Phi_F(\tfrac12+iz)=\int_\R F(t)e^{izt}\,dt=\hat F(-z),
\tag{11.2}
\end{equation}
an entire function of exponential type at most $T$ if $\supp F\subseteq[-T,T]$.
Then $\Phi_F(1-\rho)=\Phi_F(\tfrac12-iz)=\mathfrak F_F(-z)$, so
\begin{equation}
  \sum_\rho\Phi_F(\rho)\Phi_F(1-\rho)
  =\sum_{z\in\mathcal Z}\mathfrak F_F(z)\,\mathfrak F_F(-z),
\tag{11.3}
\end{equation}
absolutely convergent by Remark~\ref{rem:abs}.  If $F$ is real then
$\overline{\mathfrak F_F(\bar z)}=\mathfrak F_F(-z)$; if in addition $F$ is
even then $\mathfrak F_F$ is even and the summand is $\mathfrak F_F(z)^2$,
while if $F$ is odd then $\mathfrak F_F$ is odd and the summand is
$-\mathfrak F_F(z)^2$.

\subsection{Imposing the primitivity conditions}

\begin{lemma}
\label{lem:Pi}
Let $\Pi:=-\partial_t^2+\tfrac14$.  For $G\in C_c^\infty(\R)$,
\[
  \Phi_{\Pi G}(s)=s(1-s)\Phi_G(s),
  \qquad\text{equivalently}\qquad
  \mathfrak F_{\Pi G}(z)=(\tfrac14+z^2)\,\mathfrak F_G(z),
\]
and on $s=\tfrac12+i\tau$ with $\tau$ real the factor is $\tau^2+\tfrac14>0$.
\end{lemma}

\begin{proof}
Two integrations by parts, the boundary terms vanishing by compact support,
give $\int G''(t)e^{(s-\frac12)t}dt=(s-\tfrac12)^2\Phi_G(s)$, whence
$\Phi_{\Pi G}=\bigl(\tfrac14-(s-\tfrac12)^2\bigr)\Phi_G=s(1-s)\Phi_G$.  At
$s=\tfrac12+iz$ one has $s(1-s)=(\tfrac12+iz)(\tfrac12-iz)=\tfrac14+z^2$.
\end{proof}

\begin{corollary}
\label{cor:Pi}
$\Pi$ maps $C_c^\infty([-T,T])$ into the primitive class
$\{F:\Phi_F(0)=\Phi_F(1)=0\}$, preserves support, reality and parity, and is
injective.  Moreover $\tfrac14+z^2\ne0$ at every point of $\mathcal Z$.
\end{corollary}

\begin{proof}
$\Phi_{\Pi G}(0)=\Phi_{\Pi G}(1)=0$ by Lemma~\ref{lem:Pi}, which by
Theorem~\ref{thm:polarpacket}(ii) is primitivity.  $\Pi$ is a real,
constant-coefficient differential operator of even order, so it does not
enlarge supports, preserves reality, and commutes with $t\mapsto-t$.  For
injectivity: $\Phi_{\Pi G}\equiv0$ forces $\Phi_G\equiv0$ off $\{0,1\}$, hence
everywhere since $\Phi_G$ is entire, hence $G=0$; this is analyticity, not a
continuity argument.  Finally $\tfrac14+z^2=0$ only for $z=\pm i/2$, i.e.\ for
$s\in\{0,1\}$, which are not zeros of $\xi$ by (E2).
\end{proof}

\subsection{The construction}

\begin{theorem}[Theorem F]
\label{thm:sufficiency}
Suppose the completed zeta function $\xi$ has a zero off the line
$\Rel s=\tfrac12$.  Then there exist nonzero real $F_+\in C_c^\infty(\R)$ even
and $F_-\in C_c^\infty(\R)$ odd, both primitive
($\Phi_{F_\pm}(0)=\Phi_{F_\pm}(1)=0$), with
\begin{equation}
  \sum_\rho\Phi_{F_+}(\rho)\Phi_{F_+}(1-\rho)<0,
  \qquad
  \sum_\rho\Phi_{F_-}(\rho)\Phi_{F_-}(1-\rho)<0 .
\tag{11.4}
\end{equation}
Equivalently, by Theorem~\emph{\ref{thm:comparison}}, there are primitive
radial classes with $I_\infty(\Gamma_{F_\pm},\Gamma_{F_\pm})>0$.
\end{theorem}

\begin{proof}
\emph{Step 1: the quartet of a zero.}  Let $\rho_0$ be a zero of $\xi$ with
$\Rel\rho_0\ne\tfrac12$ and let $z_0=x_0+iy_0:=(\rho_0-\tfrac12)/i$, so
$x_0=\Iml\rho_0$ and $y_0=\tfrac12-\Rel\rho_0$.  Then $y_0\ne0$ by hypothesis,
and $x_0\ne0$ because $\zeta$ has no real zero in $(0,1)$, which is the last
clause of (E2); this is the only place in the paper where that clause is used.
Hence
\begin{equation}
  \Iml(z_0^2)=2x_0y_0\ne0 ,
\tag{11.5}
\end{equation}
and the four points of the quartet
$\mathcal O_0:=\{z_0,-z_0,\bar z_0,-\bar z_0\}$ are distinct and all lie in
$\mathcal Z$, with a common multiplicity $m_0\ge1$ by the two symmetries of the
divisor.

\emph{Step 2: a damping factor.}  Choose a nonzero real even
$\varphi\in C_c^\infty(\R)$ with $D:=\hat\varphi$ satisfying $D(z_0)\ne0$; such
a $\varphi$ exists, since for an even approximate identity concentrated at the
origin one has $\hat\varphi(z_0)\to1$.  The function $D$ is entire of finite
exponential type, real and even on $\R$, has real Taylor coefficients, and
satisfies, for every $M\ge0$,
\begin{equation}
  \sup_{|\Iml z|\le1/2}\ (1+|\Rel z|)^M|D(z)|<\infty ,
\tag{11.6}
\end{equation}
by repeated integration by parts.  In particular (11.6) with $M=1$ gives
$|D(z)|\le C_1/(1+|\Rel z|)$ in the strip $|\Iml z|\le\tfrac12$, so that
$|D(z)|\to0$ as $|\Rel z|\to\infty$ there; fix $R>|x_0|+1$ such that no point
of $\mathcal Z$ has $|\Rel z|=R$ and such that
\begin{equation}
  \sup_{\substack{|\Rel z|\ge R\\ |\Iml z|\le1/2}}
     \frac{|D(z)|}{|D(z_0)|}\le\frac12 .
\tag{11.7}
\end{equation}
Note that $D$ is fixed first and $R$ is chosen afterwards.

\emph{Step 3: the other zeros of bounded height.}  The set
$\mathcal Z_R:=\{z\in\mathcal Z:|\Rel z|<R\}$ is finite, by (E3d).  Delete from
it every occurrence of the quartet $\mathcal O_0$, with its multiplicity $m_0$
(multiset subtraction), and put
\[
  P(z):=\prod_{w\in\mathcal Z_R\setminus\mathcal O_0}(z-w),
  \qquad
  Q(z):=P(z)P(-z),
  \qquad
  E(z):=(\tfrac14+z^2)\,Q(z).
\]
The remaining divisor is invariant under $w\mapsto\bar w$, so $P$ has real
coefficients; hence $Q$ has real coefficients and is even, and so is $E$.
Since $\mathcal Z_R\setminus\mathcal O_0$ is invariant under $w\mapsto-w$, the
polynomial $Q$ vanishes at every point of $\mathcal Z_R\setminus\mathcal O_0$,
to at least its multiplicity.  Finally $E(z_0)\ne0$: no point of
$\mathcal O_0$ belongs to $\mathcal Z_R\setminus\mathcal O_0$, so
$P(z_0)\ne0$ and $P(-z_0)\ne0$, and $\tfrac14+z_0^2\ne0$ by
Corollary~\ref{cor:Pi}.

\emph{Step 4: prescribing the value at $z_0$.}  For a
positive integer $N$ consider
\begin{equation}
  \mathfrak F_{N,+}(z):=E(z)\,(a+bz^2)\,D(z)^N,
  \qquad
  \mathfrak F_{N,-}(z):=iz\,E(z)\,(c+dz^2)\,D(z)^N,
\tag{11.8}
\end{equation}
with $a,b,c,d\in\R$.  Prescribing the value at $z_0$ is the real-linear map
\begin{equation}
  \R^2\to\C,\qquad (a,b)\longmapsto E(z_0)D(z_0)^N\,(a+b\,z_0^2),
\tag{11.9}
\end{equation}
which is an isomorphism of real vector spaces if and only if $1$ and $z_0^2$
are linearly independent over $\R$ in $\C$, i.e.\ if and only if
$z_0^2\notin\R$ --- which is (11.5).  The scalar $E(z_0)D(z_0)^N$ is nonzero by
Steps 2 and 3.  We may therefore choose real $a_N,b_N$ with
$\mathfrak F_{N,+}(z_0)=i$, and real $c_N,d_N$ with
$\mathfrak F_{N,-}(z_0)=1$; and since the inverse of the fixed isomorphism
$(a,b)\mapsto a+bz_0^2$ does not depend on $N$, there is $C_R$ independent of
$N$ with
\begin{equation}
  |a_N|+|b_N|+|c_N|+|d_N|\le C_R\,|D(z_0)|^{-N}.
\tag{11.10}
\end{equation}
The coefficients $a_N,b_N,c_N,d_N$ must be real; this is why the value at $z_0$
is prescribed through the real-linear map (11.9).

\emph{Step 5: the test functions.}  The function $\mathfrak F_{N,+}$ is entire
of finite exponential type, has real Taylor coefficients and is even; hence
$\overline{\mathfrak F_{N,+}(\bar z)}=\mathfrak F_{N,+}(z)
=\mathfrak F_{N,+}(-z)$.  By (11.6) it is rapidly decreasing on $\R$.  By the
Paley--Wiener theorem (E6) there is $F_{N,+}\in C_c^\infty(\R)$, real and even,
with $\hat F_{N,+}(z)=\mathfrak F_{N,+}(-z)$, that is, with
$\mathfrak F_{F_{N,+}}=\mathfrak F_{N,+}$ in the notation (11.2); its support
is contained in a fixed interval $[-T_N,T_N]$ determined by the exponential
types.  Similarly, writing $\mathfrak F_{N,-}=iR_N$ with
$R_N(z)=zE(z)(c_N+d_Nz^2)D(z)^N$ a real-coefficient odd entire function, one
gets $\overline{\mathfrak F_{N,-}(\bar z)}=-iR_N(z)=-\mathfrak F_{N,-}(z)
=\mathfrak F_{N,-}(-z)$, so there is $F_{N,-}\in C_c^\infty(\R)$, real and odd,
with $\mathfrak F_{F_{N,-}}=\mathfrak F_{N,-}$.  Both are nonzero, since their
transforms do not vanish at $z_0$.  Both are \emph{primitive}: by construction
$\mathfrak F_{N,\pm}$ contains the factor $\tfrac14+z^2=s(1-s)$, which vanishes
at $z=\pm i/2$, i.e.\ at $s=0$ and $s=1$, so
$\Phi_{F_{N,\pm}}(0)=\Phi_{F_{N,\pm}}(1)=0$.  (Equivalently, $F_{N,\pm}=\Pi
G_{N,\pm}$ for a $G$ with $\mathfrak F_{G}=\mathfrak F_{N,\pm}/(\tfrac14+z^2)$,
which is again of Paley--Wiener class; Corollary~\ref{cor:Pi}.)

\emph{Step 6: the quartet $\mathcal O_0$ contributes $-4m_0$.}  By (11.3), the
contribution of $\mathcal O_0$ to $\sum_{z}\mathfrak F(z)\mathfrak F(-z)$ is,
in the even case,
\[
  m_0\!\!\sum_{z\in\mathcal O_0}\!\!\mathfrak F_{N,+}(z)^2
  =m_0\bigl[2\mathfrak F_{N,+}(z_0)^2+2\overline{\mathfrak F_{N,+}(z_0)^2}\bigr]
  =4m_0\,\Rel\bigl(\mathfrak F_{N,+}(z_0)^2\bigr)=4m_0\Rel(i^2)=-4m_0 ,
\]
using evenness and $\mathfrak F_{N,+}(\bar z)=\overline{\mathfrak F_{N,+}(z)}$;
and in the odd case the summand is $-\mathfrak F_{N,-}(z)^2$, so the
contribution is $-4m_0\Rel(1^2)=-4m_0$.

\emph{Step 7: the other zeros with $|\Rel z|<R$ contribute $0$.}  If
$w\in\mathcal Z_R\setminus\mathcal O_0$ then $Q(w)=0$ by Step 3, hence
$\mathfrak F_{N,\pm}(w)=0$ and the corresponding summand of (11.3) vanishes;
repeated occurrences vanish likewise.

\emph{Step 8: the remaining zeros.}  Let $D_0$ be a fixed integer exceeding
the degrees of the polynomial factors appearing in (11.8), namely
$\deg E+3$.  From (11.10), for
$|\Iml z|\le\tfrac12$,
\begin{equation}
  |\mathfrak F_{N,\pm}(z)|\le C_R(1+|z|)^{D_0}
     \Bigl(\frac{|D(z)|}{|D(z_0)|}\Bigr)^{N}.
\tag{11.11}
\end{equation}
Fix any integer $N_0\ge1$.  By the strip decay (11.6) with $M$ large together
with the zero count $\#\{z\in\mathcal Z:|\Rel z|\le X\}=O(X\log X)$ of (E3d),
\begin{equation}
  \Sigma_0:=\sum_{\substack{z\in\mathcal Z\\|\Rel z|\ge R}}(1+|z|)^{2D_0}
     \Bigl(\frac{|D(z)|}{|D(z_0)|}\Bigr)^{2N_0}<\infty .
\tag{11.12}
\end{equation}
The sum $\Sigma_0$ is indeed finite: the factor $(|D(z)|/|D(z_0)|)^{2N_0}$ decays
faster than any power of $|\Rel z|$ in the strip, since $N_0\ge1$ and (11.6)
holds for every $M$, while the number of zeros with $|\Rel z|\in[X,X+1]$ is
$O(\log X)$.
For $N\ge N_0$, note that if $|\Rel z|\ge R$ and $|\Iml z|\le\tfrac12$ then the
same holds for $-z$, and that $|D(-z)|=|D(z)|$ because $\varphi$ is even, so
that (11.7) applies to both factors; hence by (11.11) and (11.7),
\[
  \sum_{\substack{z\in\mathcal Z\\|\Rel z|\ge R}}
    \bigl|\mathfrak F_{N,\pm}(z)\mathfrak F_{N,\pm}(-z)\bigr|
  \le C_R^2\,4^{-(N-N_0)}\,\Sigma_0\;\xrightarrow[N\to\infty]{}\;0 .
\]

\emph{Conclusion.}  By Steps 6--8 and the absolute convergence of (11.3),
\[
  \sum_\rho\Phi_{F_{N,\pm}}(\rho)\Phi_{F_{N,\pm}}(1-\rho)
  =-4m_0+O\bigl(4^{-(N-N_0)}\bigr)<0
\]
for all sufficiently large $N$.  Fixing one such $N$ in each parity and setting
$F_\pm:=F_{N,\pm}$ gives (11.4); both tests are nonzero and primitive by
Step 5.  The final assertion is Theorem~\ref{thm:comparison}, identity (9.5).
\end{proof}

\begin{remark}
\label{rem:sufficiency}
Read contrapositively, Theorem~\ref{thm:sufficiency} says: if
$\sum_\rho\Phi_F(\rho)\Phi_F(1-\rho)\ge0$ for every primitive real
$F\in C_c^\infty(\R)$, then every zero of $\xi$ lies on the line
$\Rel s=\tfrac12$.  It does not provide such an inequality: by
Proposition~\ref{prop:indefinite} the form $I_A$ is indefinite with infinite
index on both sides, and by \cite{P44}, Theorem 9.11, the finite-place part of
the form takes a positive value on a primitive radial class at every finite
chart.
\end{remark}

\section{Nondegeneracy on radial classes}
\label{sec:radicalsec}

\begin{definition}
\label{def:radsec}
$\mathcal N_\cap:=\{\Gamma_h\in\DAr: I_A(\Gamma_h,\Gamma_g)=0\ \text{for all }
\Gamma_g\in\DAr\}$, the radical of $I_A$ on radial classes.
\end{definition}

\begin{theorem}[Theorem G]
\label{thm:radical}
Fix a finite chart $S$ with $S_f\ne\emptyset$ and let $c=0$.  Then for
arbitrary $h,g\in C_c^\infty(\R)$ --- no support hypothesis ---
\begin{equation}
  I_A(\Gamma_h,\Gamma_g)
  =\frac1{2\pi}\int_\R\hat h(-\tau)\overline{\hat g(-\tau)}\,
     \widehat T_S(\cz,\tau)\,d\tau
  =\frac1{2\pi}\int_\R\hat h(\tau)\overline{\hat g(\tau)}\,
     \widehat T_S(\cz,\tau)\,d\tau ,
\tag{12.1}
\end{equation}
and $\mathcal N_\cap=\{0\}$.  Consequently $I_A$ is nondegenerate on $\DAr$,
which is closed under $\star$ and $\dagger$, and on which the restriction of
$M=(\lambda_-,\lambda_+)$ is still surjective onto $\BT$.
\end{theorem}

\begin{proof}
\emph{Identity \emph{(12.1)}.}  By Lemma~\ref{lem:profile}, $\widehat{\Gamma_h}$
is carried by $\cz$ and equals $\hat h(-\tau)$ there; by
Proposition~\ref{prop:radalg}, $\Gamma_h\star\Gamma_g^\dagger=\Gamma_{h*g^\vee}$
lies again in $\DAr$; and $\widehat{\Gamma_h\star\Gamma_g^\dagger}
=\widehat{\Gamma_h}\,\overline{\widehat{\Gamma_g}}$.  Apply the Plancherel form
(2.6), which by Proposition~\ref{prop:symbol}(c) holds for every element of
$\DA$ with no support hypothesis, with $d\chi|_{\cz}=d\tau/2\pi$ and $c=0$.
The second equality is the substitution $\tau\mapsto-\tau$ together with the
evenness of $\widehat T_S$ on $\cz$: on that component the finite-place symbols
are $f_p(s)+f_p(1-s)$, invariant under $\tau\mapsto-\tau$ by
Lemma~\ref{lem:localsymbols}(i) and conjugation, and
$\widehat W_\infty(0,\tau)=\log\pi-\Rel\psi(\tfrac14+\tfrac{i\tau}2)$ is even.

\emph{Independence of the constant.}  The conclusion $\mathcal N_\cap=0$ holds
for every $c\in\R$, since Proposition~\ref{prop:symbol}(b) states that
$\widehat T_S+c$ vanishes identically on no component and has discrete zero set
on each, for every real $c$.  Unlike (9.4), this theorem is therefore not tied
to $c=0$.

\emph{The radical.}  Let $\Gamma_h\in\mathcal N_\cap$.  By (12.1),
\[
  \int_\R\hat h(\tau)\overline{\hat g(\tau)}\,\widehat T_S(\cz,\tau)\,d\tau=0
  \qquad\text{for every }g\in C_c^\infty(\R).
\]
The function $m(\tau):=\hat h(\tau)\widehat T_S(\cz,\tau)$ is locally
integrable and, by Proposition~\ref{prop:symbol}(a), is a product of a Schwartz
function with a function of at most logarithmic growth, hence tempered.  The
set $\{\hat g: g\in C_c^\infty(\R)\}$ is dense in the Schwartz space (E5), so
$m=0$ almost everywhere.  By Proposition~\ref{prop:symbol}(b) the function
$\widehat T_S(\cz,\cdot)$ is real-analytic and not identically zero, so its
zero set is discrete; concretely
$\widehat T_S(\cz,0)\ge\log\pi-\psi(\tfrac14)\approx5.3722>0$ by
Proposition~\ref{prop:symbol}(d), every finite-place term at $\tau=0$ on $\cz$
being positive because $\zeta_p=p^{-1/2}$ is real and positive there.  Hence
$\hat h$ vanishes off a discrete set; being continuous, $\hat h\equiv0$, so
$h=0$ and $\Gamma_h=0$.

\emph{Surjectivity on radial classes.}  The subalgebra $\DAr$ is $\star$- and
$\dagger$-stable by Proposition~\ref{prop:radalg}, which also gives
$M(\Gamma_h)=\bigl(\int_\R h(t)e^{t/2}dt,\ \int_\R h(t)e^{-t/2}dt\bigr)$.
These two linear functionals on $C_c^\infty(\R)$ are linearly independent: if
$\alpha\int he^{t/2}\,dt+\beta\int he^{-t/2}\,dt=0$ for every
$h\in C_c^\infty(\R)$, then the continuous function
$\alpha e^{t/2}+\beta e^{-t/2}$ vanishes identically, whence $\alpha=\beta=0$,
the two exponentials being linearly independent.  Therefore
$h\mapsto M(\Gamma_h)$ is onto $\C^2=\BT$.  Proposition~\ref{prop:tate} gives
surjectivity on all of $\DA$, which is a different statement.
\end{proof}

\begin{remark}
\label{rem:three}
Theorem~\ref{thm:radical} is not a corollary of the vanishing of the radical of
$I_A$ on all of $\DA$ (\cite{P43}, Theorem 8.4): that statement pairs against
all of $\DA$, whereas $\mathcal N_\cap$ pairs only against radial tests, and a
smaller test space can only make a radical larger.  What makes the restricted
statement true is that $\DAr$ is exactly the subalgebra whose Fourier transforms
fill the Schwartz space on the single component $\cz$, where the symbol is
already nonzero.  Conversely, on the completion
$\mathcal N_A^{\mathrm{dist}}$ the radical of the scalar part of the form is
two-dimensional, namely the plane spanned by the two isotropic directions
(Proposition~\ref{prop:rulings}).  The three statements are consistent and
concern three different pairings: the radical is $0$ on $\DA$, it is $0$ on
$\DAr$, and it is the isotropic plane on the completion.  The isotropic
directions are not radial classes: they are limits of spread families whose
$L^1$-mass escapes to $|w|_S\to\infty$.
\end{remark}

\begin{remark}
\label{rem:nocolimit}
Theorem~\ref{thm:radical} concerns a fixed chart.  Its proof does not extend
to $I_\infty$ on $\DA^\infty$, because the colimit symbol is unbounded
(Remark~\ref{rem:notafunction}).
\end{remark}

\appendix

\section{Classical facts}
\label{app:classical}

The following classical facts are used in the proofs, each where it is
named.

\begin{itemize}
\item[\textbf{(E1)}] \emph{Euler product.}
$\zeta'/\zeta(s)=-\sum_{n\ge2}\Lambda(n)n^{-s}$, absolutely convergent for
$\Rel s>1$.  \emph{Used in:} step (c) of the proof of Theorem~\ref{thm:EF},
Definition~\ref{def:A}, Theorem~\ref{thm:residues}.
\emph{Source:} \cite[\S1.1]{IK}.

\item[\textbf{(E2)}] \emph{The completed zeta function.}
$\xi(s)=\tfrac12s(s-1)\pi^{-s/2}\Gamma(s/2)\zeta(s)$ is entire, satisfies
$\xi(s)=\xi(1-s)$, and its zeros are exactly the nontrivial zeros of $\zeta$
with the same multiplicities; $\zeta$ has a simple pole at $s=1$,
$\zeta(0)=-\tfrac12$, and $\xi(0)=\xi(1)=\tfrac12\ne0$.  In addition, no
nontrivial zero of $\zeta$ is real: $\zeta(\sigma)\ne0$ for
$\sigma\in(0,1)$.  \emph{Used in:} Lemma~\ref{lem:contour},
Theorem~\ref{thm:EF}, Theorem~\ref{thm:residues}; the last clause \emph{only}
in Step 1 of the proof of Theorem~\ref{thm:sufficiency}.
\emph{Source:} \cite[\S5.1]{IK}.

\item[\textbf{(E3a)}] \emph{Hadamard product.}  $\xi$ has order $1$, and
$\xi'/\xi(s)=\sum_\rho\bigl(\frac1{s-\rho}+\frac1\rho\bigr)+B$, locally
uniformly off the zeros.  \emph{Used in:} (E3c).
\emph{Source:} \cite[ch.~12]{Davenport}.

\item[\textbf{(E3b)}] \emph{Local zero count.}
$N(X+1)-N(X)=O(\log X)$.  Consequently, for each integer $j$ there is
$X_j\in[j,j+1]$ whose distance to every ordinate of a zero is at least
$c_0/\log j$: the $O(\log j)$ ordinates in $[j,j+1]$ cannot block every
subinterval of length $\asymp1/\log j$.  \emph{Used in:}
Lemma~\ref{lem:contour}.  \emph{Source:} \cite[ch.~15]{Davenport},
.

\item[\textbf{(E3c)}] \emph{Strip estimate.}
$\xi'/\xi(\sigma\pm iX_j)=O(\log^2X_j)$ uniformly for $\sigma$ in a fixed
vertical strip, along the heights $X_j$ of (E3b).  This follows from (E3a)
together with (E3b): the $O(\log X_j)$ nearby terms of the partial-fraction
expansion each contribute $O(\log X_j)$, and the remainder is $O(\log X_j)$.
\emph{Used in:} Lemma~\ref{lem:contour}.
\emph{Source:} \cite[ch.~17]{Davenport}.
Note that (E3b) and (E3c) are not consequences of the global count
(E3d): a bound on $N(X)$ says nothing about $N(X+1)-N(X)$, and the $O(\log^2)$
bound needs the Hadamard expansion in addition.

\item[\textbf{(E3d)}] \emph{Global zero count.}
$N(X)=\frac{X}{2\pi}\log\frac X{2\pi}-\frac X{2\pi}+O(\log X)$; only
$N(X)=O(X\log X)$ is used.  \emph{Used in:} Remark~\ref{rem:abs} and Steps 3
and 8 of the proof of Theorem~\ref{thm:sufficiency}.
\emph{Source:} \cite[ch.~15]{Davenport}, \cite[\S5.1]{IK}.

\item[\textbf{(E4)}] \emph{Digamma asymptotics and special values.}
$\psi(z)=O(\log|z|)$ as $|z|\to\infty$ away from the poles (Stirling);
$\psi(z)=-1/z+O(1)$ as $z\to0$; $\psi(\tfrac12)=-\gamma-2\log2$ and
$\psi(\tfrac14)=-\gamma-3\log2-\pi/2$.  \emph{Used in:}
Theorems~\ref{thm:archsymbol}, \ref{thm:archresidue},
step (b) of the proof of Theorem~\ref{thm:EF}, \ref{thm:residues}, Lemma~\ref{lem:localsymbols}.
\emph{Source:} \cite[ch.~5]{NIST}, \cite[ch.~12]{WW}.

\item[\textbf{(E5)}] \emph{Fourier analysis on $\R$.}  Fourier inversion for
$C_c^\infty$ functions; the Plancherel theorem; and the density of
$\{\hat g: g\in C_c^\infty(\R)\}$ in the Schwartz space.  \emph{Used in:}
step (c) of the proof of Theorem~\ref{thm:EF}, Lemma~\ref{lem:traceweil},
Corollary~\ref{cor:c-unique}, Theorem~\ref{thm:radical}.
\emph{Source:} \cite[ch.~7]{Hormander}.

\item[\textbf{(E6)}] \emph{Paley--Wiener.}  An entire function of exponential
type $\le T$ which is rapidly decreasing on the real axis is the Fourier
transform of a $C^\infty$ function supported in $[-T,T]$; conversely
multiplication by a polynomial does not change the exponential type.
\emph{Used in:} Step 5 of the proof of Theorem~\ref{thm:sufficiency}.
\emph{Source:} \cite[Theorem 7.3.1]{Hormander}.

\item[\textbf{(E7)}] \emph{Lerch's formula.}
$\zeta_H'(0,\alpha)=\log\Gamma(\alpha)-\tfrac12\log2\pi$ and
$\zeta_H(0,\alpha)=\tfrac12-\alpha$ for the Hurwitz zeta function.
\emph{Used in:} Proposition~\ref{prop:zetadet}.
\emph{Source:} \cite[\S13.21]{WW}, \cite[ch.~25]{NIST}.

\item[\textbf{(E8)}] \emph{Mertens.}
$\sum_{N\le x}\Lambda(N)/N=\log x+O(1)$.  \emph{Used in:}
Proposition~\ref{prop:mertens} and, through \cite{P43}, Theorem 15.2, in
Remark~\ref{rem:mertens-scope}; nowhere else, and no statement of this paper
other than those two depends on it (Remark~\ref{rem:mertens-scope}).
\emph{Source:} \cite[\S2.2]{IK}, \cite[ch.~7]{Davenport}.
\end{itemize}

The explicit formula itself (Theorem~\ref{thm:EF}) is proved in
Section~\ref{sec:EF} from (E1)--(E5); for the classical statement see
\cite{Weil1952}, \cite{Barner1981} and \cite[\S5.5]{IK}.

\begin{thebibliography}{99}

\bibitem{Barner1981} K.~Barner,
\emph{On A.~Weil's explicit formula},
J.~Reine Angew. Math. \textbf{323} (1981), 139--152.

\bibitem{Davenport} H.~Davenport,
\emph{Multiplicative Number Theory}, 3rd ed.,
Graduate Texts in Mathematics \textbf{74}, Springer, New York, 2000.

\bibitem{Hormander} L.~H\"ormander,
\emph{The Analysis of Linear Partial Differential Operators I}, 2nd ed.,
Springer, Berlin, 1990.

\bibitem{IK} H.~Iwaniec and E.~Kowalski,
\emph{Analytic Number Theory},
American Mathematical Society Colloquium Publications \textbf{53},
Providence, RI, 2004.

\bibitem{NIST} F.~W.~J.~Olver, D.~W.~Lozier, R.~F.~Boisvert and C.~W.~Clark
(eds.), \emph{NIST Handbook of Mathematical Functions},
Cambridge University Press, Cambridge, 2010.

\bibitem{Weil1952} A.~Weil,
\emph{Sur les ``formules explicites'' de la th\'eorie des nombres premiers},
Comm. S\'em. Math. Univ. Lund (1952), 252--265.

\bibitem{WW} E.~T.~Whittaker and G.~N.~Watson,
\emph{A Course of Modern Analysis}, 4th ed.,
Cambridge University Press, Cambridge, 1927.

\bibitem{P43} D.~A.~Trejo Pizzo,
\emph{A correspondence algebra with a canonical trace on the idele class
group, its Tate symbol and its numerical form}, preprint, 2026.

\bibitem{P44} D.~A.~Trejo Pizzo,
\emph{An arithmetic Frobenius family and its integral shadow on a
correspondence algebra of the idele class group}, preprint, 2026.

\end{thebibliography}

\end{document}
